\chapter{Lecteurs bidirectionnels \texorpdfstring{\(K_D\)}{K(D)} et \texorpdfstring{\(K_\Omega\)}{K(Omega)} : profils, coïncidences, tours \texorpdfstring{\(90/120\)}{90/120} et opérateurs des familles de trajectoires}
\label{chap:lectores-bidireccionales-torres-operadores}

\section{Lecteurs \(K_D\) et \(K_\Omega\) : profils dodécaphasiques, domaine des parcours et codomaine spectral}
Pour chaque parcours orienté, on définit
\[
 K_D=\left(D_{27}+D_{36},D_1,\frac{D_4-D_3}{2}\right),
 \qquad K_\Omega=(\Omega,54,P_6).
\]
Le recensement exact des 324 parcours contient quatorze profils \(K_D\), neuf profils \(K_\Omega\), quarante couples et dix-huit coïncidences, toutes de valeur \((80,54,6)\). Ces chiffres décrivent des lecteurs, non des espèces matérielles.

\section{Semi-groupe et opérateurs de fibre}
Les hauteurs globales forment
\[
 \langle90,120\rangle=\{90,120\}\cup\{30r:r\ge6\},
\]
avec l’algèbre \(\Bbbk[X_{90},Y_{120}]/(X_{90}^4-Y_{120}^3)\). Les coefficients de tour convergents raffinent le caractère et définissent des opérateurs positifs sur les treize multisections.

\section{Construction électronique, transducteur de tours et horloge de parcours}
L’électron central fournit le cas de rang un ; les autres multisections exigent l’opérateur de fibre, le couplage et la lecture spectrale. Le transducteur conserve séparément longueur temporelle, mémoire, retenue et hauteur harmonique.

\section{Registre dodécaphasique et mémoire \texorpdfstring{\(1+5+6\)}{1+5+6} du TPK}

Le sélecteur opératoire, le calendrier et le registre dodécaphasique fixent la décomposition de mémoire utilisée par les lecteurs de masse. L’identité \(1+5+6\) appartient à l’état temporel du TPK et n’est pas reconstruite à partir des sorties spectrales.

\begingroup
\let\chapter\section
\let\section\subsection
\let\subsection\subsubsection
\let\subsubsection\paragraph
\let\clearpage\relax
\let\cleardoublepage\relax
\section{Système temporel dodécaphasique orienté du TPK}
\label{sec:ciclo-dodecafasico-tpk-rev7}
\index{TPK!système dodécaphasique}\index{périodicité d'ordre 108}

\HMTClaim{HMT-R-TPK-GAMMA108}

\subsection{Opérateur de sélection de route et transport vers l'observable dimensionnelle}

La fonction ternaire de phase opératoire est
\[
M_{\rm ph}(r)=
\begin{cases}
 +1,&r\in\{1,4,7\},\\
 -1,&r\in\{2,5,8\},\\
 0,&r\in\{3,6,9\}.
\end{cases}
\]
Ses trois valeurs sont les entrées scalaires avec lesquelles
\(\operatorname{Sel}_t\) active, respectivement, le curseur additif, le
curseur multiplicatif ou le compteur neutre. L'échantillonnage stroboscopique est
\[
 \operatorname{Obs}_3(a)=(a_3,a_6,a_9,\ldots).
\]
La sélection \(\operatorname{Sel}_t[\,M_{\rm ph}(g(t))\,]\) décide quelle
évaluation opère ; le transport cardinal déplace l'état ; et
\(\operatorname{Obs}_3\) sélectionne les instants qui sont publiés.

Le transport cardinal agit par quatre translations
\[
 T_d(i,j)=(i,j)+\nu_d\pmod9,
\]
\[
 \nu_N=(-1,0),\qquad \nu_E=(0,1),\qquad
 \nu_S=(1,0),\qquad \nu_O=(0,-1).
\]
Les curseurs additif et multiplicatif conservent des positions et des orientations
séparées. La coordonnée TRIT, la fonction \(M_{\rm ph}\), l'opérateur
\(\operatorname{Sel}_t\), les quatre translations et
\(\operatorname{Obs}_3\) sont donc des objets de types distincts.

\subsection{Calendrier et registre dodécaphasiques}

Le noyau topologique de phase ne s'arrête pas à la mise à jour d'une position
sur \(\mathbb Z_9^2\). Son état temporel minimal conserve l'ordre des
événements, l'orientation, le retour et la mémoire des transitions. Pour
une orbite de \(108\) pas, on introduit la partition ordonnée
\[
 E_m=\{9(m-1)+1,\ldots,9m\},
 \qquad m=1,\ldots,12,
\]
et l'on note son support ordonné
\[
 \Gamma_{108}=E_1\sqcup\cdots\sqcup E_{12},
 \qquad 108=12\cdot9.
\]
De manière équivalente,
\[
 \Theta_{108}=\mathbb Z/108\mathbb Z,\qquad
 \beta(\theta)=
 \left[\left\lfloor\frac{\widetilde\theta}{9}\right\rfloor\right]_{12},
 \qquad
 r(\theta)=1+(\theta\bmod9).
\]
La coordonnée \(\beta\) sélectionne l'un des douze secteurs et \(r\) localise la
phase dans sa fenêtre nonadique.

\begin{proposition}[Extension centrale non scindée de l'horloge observable]
\label{prop:extension-central-108-no-escindida-rev8}
En écrivant \(C_n=\mathbb Z/n\mathbb Z\), le calendrier observable forme la
suite exacte
\[
 0\longrightarrow C_{12}
 \xrightarrow{\ \iota\ }C_{108}
 \xrightarrow{\ p\ }C_9
 \longrightarrow0,
 \qquad
 \iota([b]_{12})=[9b]_{108},\qquad
 p([\theta]_{108})=[\theta]_9.
\]
Pour la section ensembliste \(s([r]_9)=[r]_{108}\), avec
\(0\leq r<9\), la couture est donnée par le cocycle de retenue
\[
 c(r,r')=
 \left[\left\lfloor\frac{r+r'}{9}\right\rfloor\right]_{12},
 \qquad
 s(r)+s(r')=s([r+r']_9)+\iota\bigl(c(r,r')\bigr).
\]
L'extension ne se scinde pas. Pour l'action triviale,
\[
 H^2(C_9,C_{12})
 \simeq C_{12}/9C_{12}
 \simeq C_3,
\]
et \(c\) représente la classe génératrice \([1]\).
\end{proposition}

\begin{proof}
L'image de \(\iota\) est le sous-groupe des multiples de neuf dans
\(C_{108}\), qui coïncide avec \(\ker p\) ; d'où l'exactitude. L'identité
du cocycle est la division euclidienne de \(r+r'\) par neuf. Si l'extension
centrale se scindait, le groupe médian serait isomorphe à
\(C_{12}\times C_9\). Cependant,
\[
 \exp(C_{108})=108,\qquad
 \exp(C_{12}\times C_9)=\operatorname{mcm}(12,9)=36,
\]
ce qui l'empêche. L'identification cohomologique pour un groupe cyclique à
action triviale donne \(H^2(C_9,C_{12})\simeq C_{12}/9C_{12}\), et la retenue
unitaire se projette sur le générateur de ce quotient.
\end{proof}

La non-scission type le retour : neuf itérations ferment la coordonnée dans
\(C_9\) et translatent d'une unité dans \(C_{12}\) ; cent huit itérations ferment les deux
coordonnées de cette horloge. Le calendrier \(C_{108}\) est la projection
temporelle de l'état enrichi \(\mathcal X_{\rm TPK}^{\rm enr}\) :
orientation, feuillet, frontière, cylindre, généalogie et mémoire profonde
restent dans ce dernier et peuvent
distinguer des histoires ayant le même calendrier observable.

Si \(g_0(\theta)=\theta\bmod9\), alors
\[
 g_0(\theta+9)=g_0(\theta),\qquad
 \beta(\theta+9)=\beta(\theta)+1\pmod{12}.
\]
Le témoin
\[
 (g_0,r,\beta,30^\circ\beta)(49)=(4,5,5,150^\circ),
\]
\[
 (g_0,r,\beta,30^\circ\beta)(58)=(4,5,6,180^\circ)
\]
montre exactement ce qui revient et ce qui conserve la mémoire : la phase nonadique et
la position interne coïncident, tandis que le secteur dodécaphasique avance d'une unité.
Le calibre régional \(\operatorname{diag}_c\) des régions de \(\pi\) n'est
pas cette coordonnée \(\beta\). Il n'existe pas non plus ici de morphisme
\(30^\circ\beta\to h\in\langle90,120\rangle\) ; en particulier,
\(150\notin\langle90,120\rangle\). La paire \(49/58\) est un témoin de retenue,
non un sélecteur régional ou spectral.

Pour une histoire enrichie \(x\) de TPK, soient \(\tau_m\) la sous-route
ordonnée sur \(E_m\), \(\sigma_m\) son orientation, \(\kappa_m\) la retenue
qui franchit la frontière du bloc et \(\Lambda_m\) l'état accumulé du
registre. L'objet
\[
 \mathcal R_{12}
 =\bigl((E_m,\tau_m,\sigma_m,\kappa_m,\Lambda_m)\bigr)_{m=1}^{12}
\]
est le système temporel dodécaphasique orienté du TPK. Il ne constitue pas une quatrième
théorie ajoutée à APP, TRIT et TPK : c'est la forme temporelle interne avec laquelle
TPK ordonne une orbite, conserve ses retours et publie un registre
réversible. On distingue les projections canoniques
\[
 \mathcal X_{\rm TPK}^{\rm enr}
 \longrightarrow \mathcal R_{12}
 \longrightarrow \Theta_{108}.
\]
La première conserve \((\tau_m,\sigma_m,\kappa_m,\Lambda_m)\) ; la seconde
publie uniquement la suite observable et ne peut pas reconstruire à elle seule
l'état signé. Le registre enrichi est noté \(\mathcal R_{12}\) et la
rotation des secteurs est notée \(C_{12}\).

La fermeture observable du calendrier satisfait
\[
 F_{\rm obs}^{108}=I.
\]
Les positions et les orientations sont déjà restaurées après \(54\) itérations,
tandis que \(\mathcal R_{12}\) distingue la position temporelle au sein des
douze secteurs. Ici \(F_{\rm obs}\) est la permutation observable, tandis que
\(\mathcal G_9\) est l'élagage au sein de la relation d'extension et de frontière.
La monodromie est la propriété conjointe
\[
 g(K+9)=g(K),\qquad v_{K+9}\ne v_K,
\]
non le nom isolé de l'élagage. La fermeture chronologique restitue le calendrier
observable ; le retour nonadique conserve la phase et modifie le préfixe, le cylindre
et la mémoire.

On distingue également deux mots de douze composantes :
\[
 U_{12}^{\rm cat}=U_6\Vert U_6,
 \qquad
 U_{12}^{\rm sgn}=\mathcal H_{12}(K).
\]
Le premier concatène deux émissions de longueur six. Le second est une
coordonnée signée de l'état enrichi et reste lié de façon réversible à
\(K\) par la transformation dodécaphasique \(\mathcal H_{12}\). Ni leurs domaines
ni leurs fonctions ne sont interchangeables ; en particulier, le mot concaténé
ne construit pas à lui seul la coordonnée signée.

\begin{definition}[Registre entier des douze événements]
\label{def:registro-dodecafasico-rev7}
Pour une route de chiffres \(d_1,\ldots,d_{108}\), soit
\[
 \mathcal D_{\rm evt}=\{2,4,6,8\},\qquad
 \Sigma_{12}=(+,+,+,-,-,-,+,+,+,-,-,-).
\]
Son mot orienté d'événements est
\[
 e_m=\Sigma_{12}(m)\,
 \#\{t\in E_{m+1}:d_t\in\mathcal D_{\rm evt}\},
 \qquad 0\le m<12.
\]
Les deux demi-couronnes sont appariées au moyen de
\[
 p_i=e_i+e_{i+6},\qquad
 a_i=e_i-e_{i+6},\qquad 0\le i<6.
\]
On définit le comptage brut, les cinq différences par rapport à la sixième paire et
les six antisymétries par
\[
 s_C=\sum_{i=0}^{5}p_i,\qquad
 M_i=p_i-p_5\quad(0\le i<5),\qquad
 A_6=(a_0,\ldots,a_5).
\]
Le registre entier est
\[
 \mathscr L_{12}(e)
 =\bigl(s_C,M_0,\ldots,M_4,a_0,\ldots,a_5\bigr).
\]
\end{definition}

\begin{theorem}[Coordonnées exactes de la mémoire dodécaphasique]
\label{thm:memoria-1-5-6}
L'application \(e\mapsto\mathscr L_{12}(e)\) est injective sur son image.
L'inverse est donné par
\[
 p_5=\frac{s_C-\sum_{i=0}^{4}M_i}{6},\qquad
 p_i=p_5+M_i\quad(0\le i<5),
\]
\[
 e_i=\frac{p_i+a_i}{2},\qquad
 e_{i+6}=\frac{p_i-a_i}{2}.
\]
En particulier, la distribution des coordonnées est
\[
 \boxed{1+5+6=12}.
\]
\end{theorem}

\begin{proof}
Par définition,
\[
 s_C=\sum_{i=0}^{4}(p_5+M_i)+p_5
 =6p_5+\sum_{i=0}^{4}M_i.
\]
Cela détermine \(p_5\) puis tous les \(p_i\). Additionner et soustraire les
équations \(p_i=e_i+e_{i+6}\) et \(a_i=e_i-e_{i+6}\) restitue les douze
événements. La divisibilité par six et les six conditions de parité
caractérisent l'image entière ; ni une masse ni une unité n'interviennent.
\end{proof}

La division par six ne s'applique pas pendant le comptage et n'est pas un coefficient
primitif. On forme d'abord le comptage orienté brut \(s_C\) ; ensuite, on
reconstruit le registre des douze événements ; ce n'est qu'alors qu'apparaît
\(s_C/6\) dans le caractère. De manière équivalente,
\[
 W_\pm=6n_A\pm s_C
\]
sont entiers et la factorisation réticulaire possède la forme normale de Smith
\((1,12)\). Cette antériorité sera obligatoire dans la loi des masses.

Les entiers \(90\) et \(120\) réapparaissent plus loin dans six constructions
de types distincts. Pour empêcher que le nom du cardinal ne se substitue à
l'objet, on fixe dès maintenant la convention suivante :

\begin{enumerate}
 \item \(\mathcal E_{\rm dod}=(D_3,D_4,Q):\mathbb Z^{12}\to\mathbb Z^{25}\),
       extracteur entier qui reconstruit \(K\) ;
 \item \(D_{\rm raw}^{90,120}(q)\), différence diagnostique de deux séries dans
       une carte fixée ;
 \item \(\iota_{6,8}:\mathcal F_6\hookrightarrow\mathcal F_8\), inclusion
       combinatoire de cardinalités \(90\) et \(120\) ;
 \item \(K_{6\to8}=12\Delta S_{90}-\Delta S_{120}\), fonctionnelle angulaire
       définie après fixation des canaux, des séries et de la normalisation ;
 \item \(\mathfrak S=\langle90,120\rangle\), semi-groupe infini de hauteurs
       spectrales ;
 \item \(\varepsilon_{\rm res}^{\circ}\), résidu d'une identité angulaire
       exacte, non une extraction finie de l'un quelconque des cinq objets
       précédents.
\end{enumerate}

La coïncidence partielle des indices n'établit pas l'égalité des opérateurs,
des domaines, des codomaines ni de la force probante. En particulier,
\(K_{6\to8}\) utilise deux sous-queues de la hiérarchie, mais n'est pas la hiérarchie
complète ; et aucune de ces sous-queues ne forme un bloc indépendant de la loi
des masses.

\begin{remark}[Séparation canonique des objets temporels]
La notation distingue le registre orienté \(\mathcal R_{12}\), le
calendrier observable \(\Theta_{108}\), la rotation sectorielle \(C_{12}\) et
l'orbite de transport \(\Gamma_{108}\), chacun avec son domaine, son
codomaine et sa loi de composition.
\end{remark}

\endgroup

\begingroup
\let\chapter\section
\let\section\subsection
\let\subsection\subsubsection
\let\subsubsection\paragraph
\let\clearpage\relax
% Vista pública derivada de 73_rev2_electron_lectores_torres.tex: cuerpo exacto y único retítulo académico del lector temporal de 108 estados. El fragmento del handoff permanece byte-exacto.
\section{Construction complète de l'électron bêta : parcours central, signature, caractère et masse de référence}
\Needspace{7\baselineskip}
\subsection{Parcours central de l'électron et séquence d'enregistrements généalogiques}
L'électron occupe l'unique centre \((5,5)\), avec l'orientation active \(++\) :

\[
 \Gamma_e=(5,5,++),\qquad
 \tau_{12}=+++---+++---.
\]
Ses lecteurs de déplacement sont

\[
 (D_1,D_3,D_4,D_{27},D_{36})=(54,56,68,48,32).
\]
Il ne faut pas confondre :

\begin{enumerate}
\item la signature brute \((24,0,12,0,-12)\) ;
\item la signature paire CT--108 \((-12,-4,12,-6,12)\) ;
\item le vecteur de douze événements ;
\item l'origine affine \(v_e=0\) du caractère relatif.
\end{enumerate}
Il ne faut pas non plus confondre \(\Gamma_e\) avec le mot \(w_e=201101\) du lecteur de la constante d'Euler \(e\).

\Needspace{7\baselineskip}
\subsection{Sélecteur de base de l'électron et transition vers la valeur bêta de référence}
Le sélecteur interne, construit sans utiliser la masse observée, établit

\[
 -22=-\left|((\mathbb Z/9\mathbb Z)\times\mathbb F_3)\setminus\iota(R_\pi)\right|,
 \qquad
 +15=|R_\pi\times\mathbb F_3|.
\]
La lecture des lacunes fournit en outre

\[
 V_e=\begin{pmatrix}21&22\\6&7\end{pmatrix},\qquad
 \det V_e=15.
\]
Le signe négatif de 22 enregistre le déficit/la coupure ; le signe positif de 15, la rétention/l'aire orientée. La matrice est locale à l'électron et ne s'étend pas en sélecteur universel.

L'étape de base est

\[
 e_0=\frac{\sqrt3}{4}
 \left[\exp\!\left(\frac{\varphi}{\pi^2}\right)
 -22\alpha_{\rm HMT}^3\right](1+15\Delta_4),
\]
\[
 e_0=0.510998922488\ldots\ \mathrm{MeV}.
\]
Les généalogies entières sont

\[
 22=\frac{396}{18}=27-5=24-2,
 \qquad
 15=\frac{270}{18}=\binom62=\binom64.
\]
\Needspace{7\baselineskip}
\subsection{Lecteurs bidirectionnels et bêta}
Pour tout parcours,

\[
 K_D(\Gamma)=\left(D_{27}+D_{36},D_1,\frac{D_4-D_3}{2}\right),
\]
\[
 K_\Omega(\Gamma)=(\Omega_{90,120},54,P_6).
\]
Sur les 324 parcours apparaissent 14 profils \(K_D\), 9 profils \(K_\Omega\), 40 paires conjointes et 18 coïncidences. L'unique valeur commune est

\[
 (80,54,6).
\]
Ce 80 procède du lecteur bidirectionnel du parcours électronique, et non de la période 80 du mot d'Euler. Ainsi

\[
 \kappa_e=80-54\alpha_{\rm HMT}+6\Delta_4,
\]
\[
 \boxed{
 \mathfrak e_e^\beta=e_0R_{\rm act}^{\kappa_e}
 =0.5109989506900008086082571648\ldots\ \mathrm{MeV}.}
\]
C'est le résultat électronique de référence. La valeur de base ne le remplace pas. L'égalité avec la masse de l'antiparticule procède de la conjugaison du parcours et de la parité de l'opérateur de masse, et non de l'ajout d'une seconde ligne numérique.

\Needspace{7\baselineskip}
\section{Réalisation physique par couplage spectral}
\Needspace{7\baselineskip}
\subsection{Obstruction à la sélection ponctuelle équivariante d'une multisection non centrale}
Soit \(F\subset\mathcal R_{324}\) une fibre de parcours compatible avec une classe structurelle. Son espace linéaire est

\[
 \mathcal H_F=\bigoplus_{\Gamma\in F}\mathbb C e_\Gamma.
\]
Le caractère central et les lecteurs définissent des opérateurs diagonaux :

\[
 \widehat{\mathcal A}_F^{(0)}e_\Gamma
 =\chi_{\rm full}(\sigma_\Gamma,\eta_\Gamma)e_\Gamma,
\]
\[
 \widehat K_F^r e_\Gamma=\kappa_r(\Gamma)e_\Gamma,
 \qquad r\in\{D,\Omega\}.
\]
La correction two--way est

\[
 \widehat{\mathcal M}_F^{\beta,r}
 =R_{\rm act}^{\widehat K_F^r/2}
  \widehat{\mathcal M}_F^{(0)}
  R_{\rm act}^{\widehat K_F^r/2}.
\]
Lorsque \(\dim\mathcal H_F>1\), ces opérateurs ont généralement plusieurs valeurs propres. En choisir une parce qu'elle est proche d'une masse connue inverserait la causalité. Il n'est pas non plus correct d'exiger qu'un nom physique désigne une cellule unique : l'espèce non centrale réside dans une représentation de la fibre.

\Needspace{7\baselineskip}
\subsection{Foncteur spectral de la réalisation physique : noyau positif de couplage, sous-espace persistant et mesure spectrale}
Une préparation physique \(P\) fournit des données de représentation, et non une masse :

\[
 \mathfrak r(P)=
 (q,\,j,\,\epsilon_{2\pi},\,\epsilon_{4\pi},\,
  \text{statistique},\,\text{génération},\,\text{saveur},\,
  \text{orbite de couleur},\,\text{canal d'interaction}).
\]
Soit \(\mathfrak M\) un couplage compatible avec cette représentation. On définit un noyau positif

\[
 K_{P,\mathfrak M}:F\times F\longrightarrow\mathbb C
\]
qui satisfait :

\begin{enumerate}
\item \(K=K^*\) y \(K\ge0\);
\item \(K(\Gamma,\Gamma')=0\) si les parcours violent la charge, la parité, la couleur, la statistique ou
le canal de \(P\) ;

\item \(K\) commute avec les symétries que le couplage ne brise pas ;
\item \(K(C_M\Gamma,C_M\Gamma')=\overline{K(\Gamma,\Gamma')}\) par conjugaison ;
\item ses entrées sont calculées à partir de l'enregistrement généalogique, de la frontière, de la mémoire, de la retenue et de l'holonomie, sans
grandeurs externes.

\end{enumerate}
L'opérateur d'évolution couplé est

\[
 \mathcal U_{P,\mathfrak M}
 =K_{P,\mathfrak M}^{1/2}\,\mathcal U_{\rm TPK}\,
  K_{P,\mathfrak M}^{1/2}.
\]
Le sous-espace persistant s'obtient par projection spectrale sur les modes qui survivent au retour et à l'élagage :

\[
 \Pi_{P,\mathfrak M}^{\rm surv}
 =\mathbf1_{\mathcal S_{\rm surv}}(\mathcal U_{P,\mathfrak M}).
\]
La réalisation physique correcte est le foncteur spectral

\[
 \boxed{
 \mathcal S_{\rm phys}^{\rm spec}(P,\mathfrak M)
 =\left(
 \mathcal H_{P,\mathfrak M}^{\rm surv},
 \widehat{\mathcal M}_{P,\mathfrak M},
 \mu_{P,\mathfrak M}
 \right),}
\]
où

\[
 \mathcal H_{P,\mathfrak M}^{\rm surv}
 =\Pi_{P,\mathfrak M}^{\rm surv}\mathcal H_F,
\]
\[
 \widehat{\mathcal M}_{P,\mathfrak M}
 =\Pi_{P,\mathfrak M}^{\rm surv}
  \widehat{\mathcal M}_F
  \Pi_{P,\mathfrak M}^{\rm surv},
\]
et \(\mu_{P,\mathfrak M}\) est sa mesure spectrale dans l'état préparé.

Lorsque le couplage est spécifié par une représentation irréductible \(\lambda\) d'un groupe fini \(G_P\), le projecteur correspondant s'obtient sans choisir de parcours distingué :

\[
 P_{P,\lambda}
 =\frac{\dim\lambda}{|G_P|}
  \sum_{g\in G_P}\overline{\chi_\lambda(g)}\,U(g).
\]
Le résultat bidirectionnel est alors la paire comprimée

\[
 \mathbb M(P,\lambda)=
 \left(
 P_{P,\lambda}\widehat{\mathcal M}^{\beta,D}P_{P,\lambda},
 P_{P,\lambda}\widehat{\mathcal M}^{\beta,\Omega}P_{P,\lambda}
 \right).
\]
Si les deux opérateurs commutent avec la représentation irréductible, le lemme de Schur impose leur scalarité sur le bloc. S'ils ne commutent pas, le résultat correct est le spectre du bloc invariant minimal.

\Needspace{7\baselineskip}
\subsection{Naturalité du foncteur spectral sous les entrelaceurs de la dynamique TPK, du couplage et de l'opérateur de masse}
Si \(V:F\to F'\) entrelace la dynamique TPK, le noyau de couplage et les opérateurs de masse, alors

\[
 V\Pi_{P,\mathfrak M}^{\rm surv}
 =\Pi_{P',\mathfrak M'}^{\rm surv}V,
 \qquad
 V\widehat{\mathcal M}_{P,\mathfrak M}
 =\widehat{\mathcal M}_{P',\mathfrak M'}V.
\]
Par le calcul fonctionnel, la mesure spectrale est transportée par \(V\). Le résultat ne dépend ni d'une numérotation des cellules ni du choix d'une base au sein d'une orbite. C'est la condition qui doit remplacer tout sélecteur nominal non équivariant.

\Needspace{7\baselineskip}
\subsection{Conditions d'existence d'une masse scalaire sous le lecteur dimensionnel}
Une masse scalaire est déterminée dans chacun des cas suivants :

\begin{enumerate}
\item le sous-espace persistant est de rang un ;
\item une représentation irréductible fixe un unique projecteur spectral ;
\item le protocole de mesure définit un état positif normalisé
\(\omega_{P,\mathfrak M}\) sur l'algèbre des observables ;

\item une résonance définit un pôle isolé de la résolvante.
\end{enumerate}
Dans le troisième cas,

\[
 m(P,\mathfrak M)
 =\omega_{P,\mathfrak M}
  (\widehat{\mathcal M}_{P,\mathfrak M}).
\]
En l'absence de l'un de ces mécanismes, le résultat scientifique est l'opérateur, son spectre et ses multiplicités. Ce résultat n'est pas incomplet : c'est l'objet déterminé par la dynamique.

\Needspace{7\baselineskip}
\subsection{Construction finie du noyau de masse par filtres de représentation, de persistance et de couplage, suivie d'une normalisation irréductible}
Le noyau s'obtient en quatre étapes finies :

\begin{enumerate}
\item \textbf{Filtre de représentation.} On retient les parcours dont le secteur, la branche de charge,
la couleur, l'holonomie \(2\pi/4\pi\), la parité et la statistique coïncident avec \(P\).

\item \textbf{Filtre de persistance.} On utilise l'ordre partiel formé par le retour de classe,
le retour d'orbite, le retour de cellule, la séparation des préfixes, la dette de retenue et la stabilité décimale. Le résultat peut être une orbite de Pareto ; on n'impose pas l'unicité.

\item \textbf{Filtre de couplage.} On évalue l'incidence de frontière et le canal
d'interaction. Les parcours découplés reçoivent un poids nul.

\item \textbf{Normalisation.} Sur chaque composante irréductible, la trace du
noyau est normalisée à un.

\end{enumerate}
L'atlas orienté fournit toutes les données des deux premières étapes. La troisième exige que chaque réalisation physique déclare son canal ; la quatrième est canonique une fois les composantes irréductibles fixées.

\Needspace{7\baselineskip}
\section{Transducteur généalogique des tours}
\Needspace{7\baselineskip}
\subsection{Paire de mots associés à un même parcours généalogique}
Chaque \(\Gamma\in\mathcal R_{324}\) détermine simultanément un mot temporel

\[
 w_\Gamma\in\mathbb F_3^{108}
\]
et un mot dodécaphasique signé

\[
 \tau_\Gamma\in\{-1,0,+1\}^{12}.
\]
Le premier conserve les déplacements le long de CT--108 ; le second conserve l'orientation accumulée des douze événements. Les deux mots sont liés par le même enregistrement généalogique, mais ne sont pas interchangeables. Le raffinement des tours lit les deux et conserve, en outre, feuillet, retenue, enroulement et mémoire de retour.

\Needspace{7\baselineskip}
\subsection{Correspondance entre hauteur harmonique et longueur temporelle}
Écrivons

\[
 \mathfrak S=30\langle3,4\rangle
 =\{90,120\}\cup\{30r:r\ge6\}.
\]
Pour \(h=30r\), la longueur temporelle associée est

\[
 \ell_h=9r.
\]
En particulier,

\[
 90\longleftrightarrow27,\qquad
 120\longleftrightarrow36,\qquad
 360\longleftrightarrow108.
\]
Cette correspondance n'identifie pas la tour au temps : elle entrelace la hauteur de l'opérateur harmonique avec la longueur du segment de parcours qui l'engendre.

\Needspace{7\baselineskip}
\subsection{Lecteur de déplacement du cycle temporel de 108 états}
Nous définissons l'écart cyclique

\[
 D_\Gamma(n)=
 \sum_{t=0}^{107}
 \mathbf1_{\{w_{\Gamma,t}\ne
 w_{\Gamma,t+n\!\!\!\pmod{108}}\}}.
\]
Le coefficient relatif de la branche temporelle est

\[
 \boxed{\eta_{30r}^{D}(\Gamma)
 =D_\Gamma(9r)-D_{\Gamma_e}(9r).}
\]
La soustraction n'utilise pas une masse comme origine : elle utilise le parcours central comme origine affine du caractère relatif. On a

\[
 \eta_{90}^{D}+\eta_{120}^{D}
 =[D_\Gamma(27)+D_\Gamma(36)]-80,
\]
de sorte que la première couche prolonge exactement la première composante du lecteur \(K_D\).

\Needspace{7\baselineskip}
\subsection{Lecteur d'accumulation dodécaphasique}
Pour \(r\ge1\), soit

\[
 A_\Gamma(r)=
 \sum_{j=0}^{11}
 \left(
 \sum_{u=0}^{r-1}
 \tau_{\Gamma,j+u\!\!\!\pmod{12}}
 \right)^2.
\]
Le coefficient relatif de la branche orientée est

\[
 \boxed{\eta_{30r}^{\Omega}(\Gamma)
 =A_\Gamma(r)-A_{\Gamma_e}(r).}
\]
Alors

\[
 4\eta_{90}^{\Omega}-3\eta_{120}^{\Omega}
 =\Omega(\tau_\Gamma)-80,
\]
de sorte que cette couche prolonge la première composante de \(K_\Omega\). Les deux familles forment le relèvement bidirectionnel

\[
 \mathcal F_{\rm tower}^{2w}:
 \mathcal R_{324}\longrightarrow
 \mathcal E_{\mathfrak S}\times\mathcal E_{\mathfrak S},
 \qquad
 \Gamma\longmapsto
 \bigl(\eta^D(\Gamma),\eta^\Omega(\Gamma)\bigr).
\]
\Needspace{7\baselineskip}
\subsection{Détermination finie de la tour infinie}
La périodicité du mot temporel donne

\[
 D_\Gamma(9(r+12))=D_\Gamma(9r).
\]
Si

\[
 T_\Gamma=\sum_{j=0}^{11}\tau_{\Gamma,j},
 \qquad r=12q+s,\quad 0\le s<12,
\]
alors

\[
 A_\Gamma(12q+s)
 =A_\Gamma(s)+(12q^2+2qs)T_\Gamma^2.
\]
Ainsi, toute la tour est déterminée par douze valeurs de \(D_\Gamma\), douze valeurs de \(A_\Gamma\) et l'entier \(T_\Gamma^2\). L'infinité des hauteurs n'exige pas une infinité de données indépendantes.

La même formule résout le diamant généalogique de hauteur 360. Si

\[
 h(a,b)=90a+120b,
\]
alors

\[
 (a+4,b)\sim(a,b+3)
\]
puisque \(3(a+4)+4b=3a+4(b+3)\). La généalogie \((a,b)\) est conservée jusqu'à l'application de la mémoire et de la retenue ; ce n'est qu'ensuite qu'elle descend à la hauteur commune.

\Needspace{7\baselineskip}
\subsection{Convergence absolue du caractère raffiné sur les tours \(90/120\)}
On a

\[
 |\eta_h^D|\le108,\qquad
 |\eta_{30r}^{\Omega}|\le24r^2.
\]
Comme \(0<q_+<q_-<1\),

\[
 |s_{30r}|\le2q_-^{30r}.
\]
On en déduit

\[
 \sum_{h\in\mathfrak S}|\eta_h^D s_h|<\infty,
 \qquad
 \sum_{h\in\mathfrak S}|\eta_h^\Omega s_h|<\infty.
\]
Le caractère raffiné est bien défini sans tronquer la tour.

\Needspace{7\baselineskip}
\subsection{Mémoire, retenue et cocycle}
Le coefficient brut est enrichi par

\[
 \mathbb F_{\rm tower}(\Gamma)_h
 =\bigl(\eta_h^D,\eta_h^\Omega,\rho_h,
 \mathbf w_h,c_h,N_h^{\rm surv}\bigr),
\]
où la signature tronquée satisfait

\[
 \sigma_{\Gamma|\ell_h}
 =r(\rho_h)+\mathcal D\mathbf w_h,
 \qquad
 \mathcal D=\operatorname{diag}(2,6,2,2,4),
\]
et la retenue de deux résidus est

\[
 c(\rho_1,\rho_2)=
 \mathcal D^{-1}
 \bigl[r(\rho_1)+r(\rho_2)-r(\rho_1+\rho_2)\bigr].
\]
Le défaut de localité d'une composition,

\[
 \delta_h(\Gamma,\Lambda)=
 F_h(\Gamma\star\Lambda)-F_h(\Gamma)-F_h(\Lambda),
\]
obéit à l'identité de cocycle

\[
 \delta_h(\Gamma,\Lambda)
 +\delta_h(\Gamma\star\Lambda,\Xi)
 =\delta_h(\Lambda,\Xi)
 +\delta_h(\Gamma,\Lambda\star\Xi).
\]
On préserve ainsi la raison pour laquelle deux parcours de même hauteur visible peuvent transporter des histoires distinctes.

\Needspace{7\baselineskip}
\subsection{Sélection de la lecture par le couplage}
Les deux branches ne sont pas moyennées par commodité. Le couplage physique \(a\) doit déterminer une transformation naturelle

\[
 \Pi_a:
 \mathcal E_{\mathfrak S}\times\mathcal E_{\mathfrak S}
 \longrightarrow\mathcal E_{\mathfrak S}
\]
qui commute avec les troncatures et les symétries, respecte le cocycle de retenue, conserve l'origine électronique et soit compatible avec \(X_{90}^{4}=Y_{120}^{3}\). Si le couplage échange les deux lectures, l'unique clôture continue, multiplicative, symétrique et normalisée sur les scalaires positifs est la moyenne géométrique ; en termes d'exposants,

\[
 L_{\rm sym}=\frac{L_D+L_\Omega}{2}.
\]
En l'absence de cette symétrie, le résultat est la mesure spectrale conjointe de \((L_D,L_\Omega)\), et non une moyenne arbitraire.

\Needspace{7\baselineskip}
\subsection{Raffinement coinductif par orbites primitives}
Le relèvement fini précédent admet un prolongement. Soit \(\mathscr X_\Gamma\) le graphe des extensions admissibles du parcours, \(U_\Gamma\) son opérateur de transfert et \(R_\Gamma\) l'involution d'orientation. Les traces

\[
 a_n(\Gamma)=\operatorname{Tr}(R_\Gamma U_\Gamma^n)
\]
et l'inversion de Möbius

\[
 p_n(\Gamma)=\frac1n\sum_{d\mid n}\mu(d)a_{n/d}(\Gamma)
\]
séparent les orbites primitives. La différence

\[
 p_h(\Gamma)-p_h(\Gamma_e)
\]
est un raffinement ultérieur du coefficient de hauteur \(h\), et non un remplacement des deux lectures finies.

\Needspace{7\baselineskip}
\subsection{Formulation historique au moyen d'un automate d'extension}
La signature \(\mathbb Z^5\) n'épuise pas l'histoire d'un parcours. Pour chaque \(\Gamma\in\mathcal R_{324}\), on conserve :

\begin{itemize}
\item les 108 transitions de l'enregistrement généalogique ;
\item les vingt blocs et leurs préfixes ;
\item les mots \(w_6\) ;
\item les neuf comparaisons \(K\leftrightarrow K+9\) ;
\item le feuillet, la retenue et l'enroulement ;
\item les premiers retours de classe, d'orbite, de cellule et d'état cinématique ;
\item l'ambiguïté accumulée et les triades décimales stabilisées ;
\item les profils \(K_D\) et \(K_\Omega\).
\end{itemize}
Ces données permettent de construire un transducteur de tours sans consulter le résidu d'une masse.

\Needspace{7\baselineskip}
\subsection{Automate local d’extension}
Soit \(\mathscr X_\Gamma\) l’ensemble fini des états locaux compatibles avec la route. Un état conserve

\[
 x=(t,g,B,p,\varepsilon,c,w),
\]
où \(t\) est le temps discret, \(g\) la phase nonadique, \(B\) le bloc, \(p\) le préfixe, \(\varepsilon\) le feuillet, \(c\) la retenue et \(w\) l'enroulement. Il existe une arête \(x\to y\) lorsque \(y\) est une extension admissible de \(x\), respecte l'enregistrement généalogique et satisfait l'élagage \(G_9\). Le tableau des 36.864 compositions de retenue fournit la loi locale d'addition ; les tableaux de retour fixent la mémoire.

Soit \(U_\Gamma\) l’opérateur de transfert normalisé de ce graphe et \(R_\Gamma\) l’involution d’orientation. On exige

\[
 U_{C_M\Gamma}=J U_\Gamma J^{-1},\qquad
 R_{C_M\Gamma}=-J R_\Gamma J^{-1}.
\]
\Needspace{7\baselineskip}
\subsection{Dénombrements orientés de cellules primitives et construction de la signature électronique}
Pour \(n\ge1\), la trace orientée

\[
 a_n(\Gamma)=\operatorname{Tr}(R_\Gamma U_\Gamma^n)
\]
dénombre les retours fermés avec signe. L’inversion de Möbius sépare les orbites primitives :

\[
 p_n(\Gamma)
 =\frac1n\sum_{d\mid n}\mu(d)a_{n/d}(\Gamma).
\]
On prend l’électron comme origine affine des corrections relatives et on définit

\[
 \boxed{
 \eta_h(\Gamma)=p_h(\Gamma)-p_h(\Gamma_e),
 \qquad h\in\mathfrak S.}
\]
On obtient ainsi l’application

\[
 \mathcal F_{\rm tower}:
 \Gamma^{\rm enr}\longmapsto
 \eta(\Gamma)=(\eta_h(\Gamma))_{h\in\mathfrak S}.
\]
Il n’y a pas de paramètres par espèce. Le même automate, la même involution et la même inversion de Möbius agissent sur les 324 routes.

\Needspace{7\baselineskip}
\subsection{Convergence du raffinement primitif}
Soit \(N_\Gamma=\dim\mathscr X_\Gamma\). Si \(U_\Gamma\) est une contraction,

\[
 |a_n(\Gamma)|\le N_\Gamma.
\]
Le nombre de diviseurs de \(n\) croît de manière sous-linéaire, de sorte que

\[
 |p_n(\Gamma)|
 \le\frac{N_\Gamma\tau(n)}n.
\]
Comme \(0<q_+<q_-<1\),

\[
 |s_h|\le q_-^h+q_+^h\le2q_-^h.
\]
Par comparaison,

\[
 \sum_{h\in\mathfrak S}|\eta_h(\Gamma)s_h|<\infty
\]
uniformément sur toute collection finie de fibres. Le caractère raffiné est donc bien défini.

\Needspace{7\baselineskip}
\subsection{Identités de conservation du parcours électronique et unicité du point fixe central}
La construction doit satisfaire simultanément :

\begin{enumerate}
\item invariance par renumérotation des états ;
\item covariance sous conjugaison de feuillet ;
\item compatibilité avec la somme directe de composantes ;
\item égalité lors du recalcul à partir du registre généalogique et du graphe certifié ;
\item conservation de l’origine \(\eta(\Gamma_e)=0\) ;
\item insensibilité absolue à une permutation des colonnes externes de masse ;
\item détermination des queues au moyen d’une borne antérieure à la confrontation.
\end{enumerate}
La représentation exacte d’un résiduel donné démontre la complétude du codomaine, mais ne remplace pas ces sept preuves.

\Needspace{7\baselineskip}
\section{Horloge de route et normalisation absolue}
\Needspace{7\baselineskip}
\subsection{Les retours entiers ne suffisent pas à distinguer les espèces}
Dans l’atlas de 324 routes, le premier retour de cellule est 27, le retour cinématique au sein de CT--108 est 54 et le retour cinématique étendu est 216. Étant communs, ces entiers fixent l’architecture temporelle globale, mais ne peuvent pas expliquer à eux seuls les différences de masse entre espèces.

\Needspace{7\baselineskip}
\subsection{Holonomie temporelle du parcours électronique et retour avec mémoire}
L’information spécifique réside dans l’holonomie de l’opérateur d’extension. Soit \(\lambda_j(\Gamma)\) une valeur propre non triviale de \(U_\Gamma\) sur le sous-espace persistant. Nous écrivons

\[
 \lambda_j(\Gamma)=r_j(\Gamma)e^{i\theta_j(\Gamma)}.
\]
La fréquence interne associée est

\[
 \omega_j(\Gamma)
 =\frac{\theta_j(\Gamma)+2\pi w_j(\Gamma)}{108t_0},
\]
où l’entier \(w_j\) est fixé par l’enroulement enregistré, non par un choix de branche ultérieur. Pour un mode stable \(r_j=1\) ; pour une résonance \(r_j<1\) et son atténuation apporte une largeur.

Le quotient

\[
 \rho_{\omega,j}(\Gamma)
 =\frac{\omega_j(\Gamma)}{\omega_0}
\]
est adimensionnel et peut modifier les rapports de masse sans altérer la décade commune \(-34\). L’énergie et la masse du mode sont

\[
 E_j(\Gamma)=\hbar_{\rm ret}\omega_j(\Gamma)
 \chi_{\rm full}(\Gamma)R_{\rm act}^{\kappa_j(\Gamma)},
\]
\[
 m_j(\Gamma)=E_j(\Gamma)c_{\rm int}^{-2}.
\]
Cette équation réunit action, horloge, caractère et lecteur sans convertir l'un en l'autre. Son application numérique exige de déterminer \(U_\Gamma\), de fixer l'enroulement et de déclarer la carte dimensionnelle.

\Needspace{7\baselineskip}


\endgroup

\begingroup
\let\chapter\section
\let\section\subsection
\let\subsection\subsubsection
\let\subsubsection\paragraph
\let\clearpage\relax
% Vista científica exacta materializada desde una rama condicional.
% Fuente: source/demostraciones_primarias/14_05b_extension_superviviente_caracter.tex; rama: HMTIncludeRouteLedgerBlock.
% No modifica el contenido matemático de la rama de origen.
\chapter{Extension survivante, registre dodécaphasique et signature entière}
\label{chap:extension-ruta-caracter}

La trajectoire observée ne constitue pas l'objet premier de cette
construction. L'objet premier est une extension compatible à toutes les
profondeurs du transport HMT ; un chemin fini est l'une de ses lectures. Cette
distinction permet de formuler sans ambiguïté la chaîne interne qui conduit d'une
histoire enrichie au caractère d'action. Elle détermine aussi quelles
informations peuvent être oubliées sans modifier la sortie et lesquelles doivent rester dans le
registre.

Le but de cette section est de réunir en un seul argument quatre éléments qui
avaient été étudiés séparément : l'extension survivante, la mémoire
intrinsèque, le registre orienté et la signature entière. Le résultat terminal de
cette section est le caractère formel
\[
 \Gamma_{108}^{\rm enr}
 \xrightarrow{\ L_{\rm tick}\ }
 \ZZ^5
 \xrightarrow{\ \chi_{\rm form}\ }
 \operatorname{Fun}(\mathcal P,\RR_{>0}),
\]
antérieur tant à la spécialisation des coefficients qu'à toute unité
physique. La première application extrait la signature du registre lu ; la seconde la
convertit en une fonction positive sur l'espace des paramètres
\(\mathcal P\). La spécialisation numérique ne s'effectue qu'après
la construction des coefficients globaux dans la section correspondante. La
réalisation métrologique relève d'une carte dimensionnelle encore ultérieure.
L'application \(L_{\rm tick}\) est le lecteur entier d'itération et de retour du
registre chronologique enrichi.

\section{Composition opératoire de la prolongation et cochaîne torsionnelle dans la génération du registre de chemin}

La loi des masses commence avant la formule exponentielle. APP apporte un unique
support de chemins et ses deux feuillets, somme et produit. Le TRIT oriente chaque
itération élémentaire.
Si
\[
 r_t=1+((t-1)\bmod9),\qquad m_t:=M_{\rm ph}(r_t),
\]
le sélecteur de phase \(M_{\rm ph}\) active
\[
 \operatorname{op}_t=
 \begin{cases}
  \Sigma,&m_t=+1,\\
  \Pi,&m_t=-1,\\
  \varnothing\quad(\text{\rm rétention ou compteur neutre}),&m_t=0.
 \end{cases}
\]
La lecture temporelle peut interpréter le troisième cas comme seuil ou présent,
mais ne le transforme pas en un troisième feuillet arithmétique. Ce n'est pas APP qui alterne
ses feuillets. TPK transporte
conjointement position, orientation, feuillet, phase, bloc, retenue, préfixe,
cylindre, généalogie et mémoire. Le registre chronologique enrichi note
chaque transition ---position,
feuillet, lecture, retenue, phase, bloc et événement---. La signature locale de frontière
est ensuite dérivée au moyen de \(\operatorname{Sig}_q(h,c,\lambda)\), et la signature
globale de masse est extraite du registre complet ; aucune n'est reconstruite
rétrospectivement à partir d'une masse.

L'entonnoir fini qui alimente cette histoire conserve quatre types distincts :
\[
 9^2\!\cdot9^2\!\cdot4\!\cdot4=104\,976
 \longrightarrow 468\ \text{émissions }U_6
 \longrightarrow 243\ \text{mots }w_6
 \hookrightarrow \mathbb F_3^6,
 \qquad |\mathbb F_3^6|=729.
\]
Le premier chiffre compte les conditions initiales du moteur à deux curseurs ; 468
compte les émissions décimales et 243 les mots ternaires visibles. Les 324 germes
de la projection chronologique \(\Pi_{108}^{\rm cal}\) à un seul curseur
ne remplacent pas non plus les 104\,976
conditions de l'état complet. Chaque flèche change donc de domaine et
d'opération. L'égalité numérique entre l'espace ambiant de \(729\) mots et les
\(729\) suffixes d'une étape de filtration n'est reliée que par le codage
ternaire déclaré ; elle n'identifie pas les deux objets.

Le relèvement conserve également son ordre causal :
\[
 w_6\xrightarrow{\,L_0\,}w_{12}
 \xrightarrow{\,L_0\,}w_{18}
 \xrightarrow{\,L_1\,}w_{24}
 \xrightarrow{\,L_1\,}w_{30}
 \longrightarrow R_{36}\longrightarrow\mathcal G_9
 \longrightarrow
 \text{section survivante}.
\]
La frontière \(R_{36}\) ouvre la prolongation. À chaque étape de filtration,
\(\mathcal G_9\) épuise
les 729 sous-cylindres et conserve l'unique compatible avec tous les horizons ;
l'épisode spéculaire profond \(3\to1\) est un autre élagage et ne doit pas être confondu avec
cette partition ordinaire \(729\to1\). Le retour nonadique conserve la phase,
mais change le bloc, le cylindre et la mémoire. Ni chiffres cibles ni masses n'interviennent
dans cette sélection. À la frontière dimensionnelle, une particule est une section
persistante enrichie de l'extension ; le chemin est son histoire observable.
Son nom physique se lit ensuite et ne décide pas quelle branche survit.

Avant d'ouvrir le registre chronologique dimensionnel, la même généalogie conserve une
fenêtre publiée de vingt sextuors par canal et admet la première lecture
archimédienne de douze triades décimales. Ces deux témoins appartiennent à la
prolongation de l'état : ils ne sont pas les douze événements du registre de masse et ne les
remplacent pas. Les événements ne se construisent qu'ensuite, en lisant le chemin sur le
système dodécaphasique orienté.

La généalogie dodécaphasique convertit ensuite les douze événements orientés en un
registre et seulement alors en la signature
\[
 \gamma_x\longmapsto\mathcal L(\gamma_x),\qquad
 \sigma_{\rm reg}:=L_{\rm tick}(\gamma_x)
 =(n_A,s_C,k_{\Dfour},b_{120},b_{\zeta}),
\]
\[
 \sigma_{\rm reg}\longmapsto
 \chi_{\rm form}(\sigma_{\rm reg};\,\cdot).
\]
Les coordonnées publiques \(b_{120}:=\nu_{120}^{\rm ev}\) et
\(b_{\zeta}:=\nu_{270}\) sont des sorties du registre chronologique, non des corrections
métrologiques introduites à la fin. Les coefficients ultérieurs du
semi-groupe de tours conservent une autre provenance :
\[
 N_{120}=b_{120}+\eta_{120},\qquad
 b_{\zeta}\zeta_{270}\not\equiv\eta_{270}s_{270},
\]
avec \(\zeta_{270}=135\Dfour^2\) et
\(s_{270}=q_-^{270}-q_+^{270}\). Le raffinement des tours s'ajoute après
le caractère central sans effacer la provenance d'aucune coordonnée.

Trois sorties latérales conservent aussi leur type. Le secteur nul quitte la
branche positive au moyen de son lecteur et de sa carte nuls ; il n'est pas la limite
\(m\to0\) d'un caractère positif. Une trajectoire virtuelle est une extension
finie avec obstruction terminale et conserve le registre chronologique et la retenue sans acquérir pour
cela une masse imaginaire. Enfin, l'ombre de Witt accompagne chaque étape de filtration et
préserve la généalogie exceptionnelle, mais ne sélectionne ni le chemin ni les
coefficients de masse.

\begin{figure}[p]
\centering
\begin{tikzpicture}[
  font=\small,
  node distance=3.5mm,
  box/.style={draw,rounded corners=1.5mm,align=center,text width=12.8cm,
    inner xsep=6pt,inner ysep=4pt},
  hmt/.style={box,draw=hmtblue,fill=hmtlightblue},
  md/.style={box,draw=hmtviolet,fill=hmtlightviolet},
  act/.style={box,draw=hmtgreen,fill=hmtlightgreen},
  result/.style={box,draw=hmtgold,fill=hmtlightgold},
  arr/.style={-{Stealth[length=2.2mm]},line width=.7pt}]

  \node[hmt] (app) {\textbf{APP} : tore arithmétique nonadique, chemins et feuillets additif et multiplicatif. La spécification complète contient \(104\,976\) graines à double curseur.};
  \node[hmt,below=of app] (trit) {\textbf{TRIT} : orientation locale de chaque transition ; le relèvement conserve le résidu, le quotient et la retenue.};
  \node[hmt,below=of trit] (tpk) {\textbf{TPK} : transport de phase avec mémoire ; \(104\,976\to468\,U_6\to243\,w_6\subset\mathbb F_3^6\).};
  \node[hmt,below=of tpk] (surv) {Filtration coinductive et monodromie \(\mathcal G_9\) : trajectoire survivante à généalogie conservée.};
  \node[md,below=of surv] (route) {La trajectoire de particule produit, sans consulter la masse, un enregistrement dodécaphasique et la signature \(\sigma=(n_A,s_C,k_{\Delta_4},b_{120},b_\zeta)\in\mathbb Z^5\).};
  \node[md,below=of route] (chi) {Le caractère central \(\chi_0(\sigma)\) et le raffinement harmonique \(\eta\in\mathbb Z^{\langle90,120\rangle}\) construisent \(\chi_{\rm full}(\sigma,\eta)\in\mathbb R_{>0}\).};
  \node[act,below=of chi] (action) {Branche dimensionnelle indépendante : \(\Sigma_H=\varphi\alpha_{\rm HMT}^{16}\to\hbar_{\rm int}\), \(\omega_0=2\pi/(108t_0)\), \(m_0=\hbar_{\rm int}\omega_0c_{\rm int}^{-2}\).};
  \node[result,below=of action] (mass) {Confluence typée : \(m_X^{\rm int}=m_0\,\chi_{\rm full}(\sigma_X,\eta_X)\). L’échelle commune fixe la normalisation absolue et s’annule dans \(m_X/m_Y\).};
  \node[result,below=of mass] (compare) {La carte d’unités exprime la grandeur. CODATA et PDG ne sont consultés que dans la dernière colonne de comparaison.};

  \foreach \a/\b in {app/trit,trit/tpk,tpk/surv,surv/route,route/chi,
    chi/action,action/mass,mass/compare}{\draw[arr] (\a)--(\b);}
\end{tikzpicture}

\caption{Dérivation d'une grandeur de masse à partir d'une trajectoire HMT.
Le caractère apporte un rapport adimensionnel ; l'échelle d'action et l'horloge
apportent la section dimensionnelle. La référence externe n'intervient qu'après
l'application de la carte d'unités.}
\label{fig:derivacion-masa-rev9}
\end{figure}

\section{Projection d'une extension survivante sur son chemin observable}

Soit \((\mathcal S_m,\rho_{m+1,m})_{m\geq0}\) le système d'états HMT
survivants à la profondeur \(m\), avec des applications de troncature
\(\rho_{m+1,m}:\mathcal S_{m+1}\to\mathcal S_m\). Nous définissons
\[
 \Omega_{\HMT}^{\rm surv}
 :=\varprojlim_m\mathcal S_m
 =\left\{(x_m)_m:\rho_{m+1,m}(x_{m+1})=x_m\right\}.
\]
Un élément \(x\in\Omega_{\HMT}^{\rm surv}\) conserve simultanément la
compatibilité de toutes ses troncatures. Une carte de lecture de profondeur
108 produit
\[
 \operatorname{sh}_{108}(x)=\gamma_x\in\Gamma_{108}^{\rm enr}.
\]
La notation \(\operatorname{sh}\) rappelle que \(\gamma_x\) est une ombre
observable : elle enregistre une présentation séquentielle de l'état, mais l'état ne
se définit pas en choisissant rétrospectivement ce chemin.

La présentation enrichie conserve, outre le mot visible, la phase,
le feuillet, l'orientation, le bord, la mémoire antifeuillet et les prédécesseurs nécessaires pour
évaluer les cocycles. Deux mots ayant les mêmes cinq entiers finaux peuvent
être des présentations distinctes parce que ces entiers sont calculés après
l'application des signes et des corrélations.

\begin{definition}[Mémoire intrinsèque et registre lu]
Une extension dodécaphasique enrichie est une classe compatible de
troncatures. Nous notons \(\mathscr M(x)\) la mémoire intrinsèque que ces
troncatures conservent. La carte
\(\gamma_x=\operatorname{sh}_{108}(x)\) étant fixée, son registre lu est
\[
 \mathcal L(\gamma_x)
 =\bigl(n_A,e_0,\ldots,e_{11},T_\zeta,C_\zeta\bigr),
\]
où \(n_A\) est le dénombrement du support d'action déclaré, les \(e_m\) sont
les douze événements orientés et \(T_\zeta,C_\zeta\) sont la mémoire torsionnelle
totale et sa restriction à la couronne. La compatibilité exige que ces données
soient préservées par toute troncature contenant les pas employés dans leur
définition.
\end{definition}

L'extension contient déjà la mémoire que le chemin présente ; le chemin ne la crée pas.
Relativement à la carte déclarée, la composition
\[
x\longmapsto\gamma_x
\longmapsto\mathcal L(\gamma_x)
\longmapsto L_{\rm tick}(\gamma_x)
\]
est une application. L'invariance par rapport à une autre carte de présentation
exigerait en outre de démontrer la compatibilité des applications
de transition correspondantes. En ce sens précis, le registre ne s'ajoute pas à une
particule déjà construite : c'est une lecture de la mémoire de l'extension.

\section{Décomposition entière de \(12\) événements dans les secteurs \(1+5+6\)}

Soit \(d_1,\ldots,d_{108}\in\{1,\ldots,9\}\) le mot d'une présentation
enrichie. Dans la convention de parité utilisée par l'extracteur, le support
d'action est
\[
 \mathcal D_{\rm act}=\{1,3,5,7\},
 \qquad
 n_A=\#\{t:d_t\in\mathcal D_{\rm act}\}.
\]
Le choix de ce support est une partie déclarée de la carte ; il ne doit pas
être confondu avec la classification TRIT des pas. Pour les événements pairs, on
utilise
\[
 \mathcal D_{\rm evt}=\{2,4,6,8\},
 \qquad
 \Sigma_{12}=(+,+,+,-,-,-,+,+,+,-,-,-).
\]
L'événement orienté du bloc \(m\in\{0,\ldots,11\}\) est
\[
 e_m=\Sigma_{12}(m)
 \#\{t\in\{9m+1,\ldots,9m+9\}:d_t\in\mathcal D_{\rm evt}\}.
\]

Apparions les deux demi-couronnes au moyen de
\[
 p_i=e_i+e_{i+6},
 \qquad
 a_i=e_i-e_{i+6},
 \qquad 0\leq i<6.
\]
On définit
\[
 s_C=\sum_{i=0}^{5}p_i,
 \qquad
 M_i=p_i-p_5\quad(0\leq i<5),
\]
et nous écrivons \(M_5=(M_0,\ldots,M_4)\),
\(A_6=(a_0,\ldots,a_5)\).

Le \cref{thm:memoria-1-5-6}, dans la résidence primaire du système
dodécaphasique de TPK, démontre que
\[
 (e_0,\ldots,e_{11})\longmapsto(s_C,M_5,A_6)
\]
est injectif sur son image et fournit l'inverse
\[
 p_5=\frac{s_C-\sum_{i=0}^{4}M_i}{6},\qquad
 p_i=p_5+M_i,
\]
\[
 e_i=\frac{p_i+a_i}{2},\qquad
 e_{i+6}=\frac{p_i-a_i}{2}.
\]
On ne récupère ici que ces coordonnées pour construire le caractère ; la
preuve et l'identité \(1+5+6=12\) ne sont pas dupliquées.

La donnée \(s_C\) est celle qui entre dans le caractère angulaire. Les onze coordonnées
\((M_5,A_6)\) conservent le placement de ce total. Les éliminer avant de
composer des histoires confond des événements différents susceptibles de produire
des corrélations distinctes.

\section{Cochaîne torsionnelle au niveau du pas}

Le dénombrement des événements ne détermine pas à lui seul la contribution torsionnelle. Pour
chaque occurrence de neuf, nous conservons le chiffre qui la précède et définissons
\[
 j_2(d)=
 \begin{cases}
 0,&d=9,\\
 +1,&d\in\{2,4,6,8\},\\
 -1,&d\in\{1,3,5,7\}.
 \end{cases}
\]
Avec des indices cycliques dans le mot, soit
\[
 \zeta_t=\mathbf 1_{\{d_t=9\}}j_2(d_{t-1}),
\]
\[
 T_\zeta=\sum_{t=1}^{108}\zeta_t,
 \qquad
 C_\zeta=\sum_{t\in\{27,54,81,108\}}\zeta_t.
\]
La coordonnée forte qui accompagne \(\Dfour\) est
\[
 \boxed{k_{\Dfour}=T_\zeta-2C_\zeta.}
\]
La formule sépare l'entraînement accumulé de sa correction orientée dans les
quatre clôtures de couronne. Le recensement exhaustif montre que ni
\(T_\zeta\) ni \(C_\zeta\), séparément, ne séparent toutes les fibres du recensement ;
le couple ordonné le fait dans les deux conventions typées qui sont confrontées.

Les deux coordonnées de tour restantes sont calculées à partir du registre
d'événements. Soit
\[
 \tau_m=\operatorname{sgn}(e_m).
\]
Pour les quatre classes de pas quatre,
\[
 I_a=\{a,a+4,a+8\}\pmod{12},
 \qquad 0\leq a<4,
\]
nous définissons le caractère pondéré par les événements
\[
 \nu_{120}^{\rm ev}
 =\sum_{a=0}^{3}\operatorname{sgn}
 \left(\sum_{m\in I_a}e_m\right),
\]
et l'autocorrélation de pas trois
\[
 \nu_{270}=\sum_{m=0}^{11}\tau_m\tau_{(m+3)\bmod12}.
\]
La variante
\[
 \nu_{120}^{\rm sgn}
 =\sum_{a=0}^{3}\operatorname{sgn}
 \left(\sum_{m\in I_a}\tau_m\right)
\]
est une autre observable. Elles ne sont pas identifiées tacitement : elles diffèrent sur
25 des 81 chemins de l'atlas.

\begin{definition}[Extracteur enrichi de signature]
La convention d'événements étant fixée, on définit
\[
 L_{\rm tick}(\gamma)=
 \bigl(n_A,s_C,k_{\Dfour},\nu_{120}^{\rm ev},\nu_{270}\bigr)
 \in\ZZ^5.
\]
L'indice indique que la coordonnée torsionnelle est calculée avant d'oublier
les prédécesseurs des pas nuls.
\end{definition}

Cette application est déterministe et ne consulte ni masses, ni unités, ni étiquettes
de particules. On n'affirme pas qu'elle soit linéaire par rapport à une concaténation déjà
projetée : les signes et les autocorrélations imposent de composer d'abord les
registres enrichis.

\endgroup
\begingroup
\let\chapter\section
\let\section\subsection
\let\subsection\subsubsection
\let\subsubsection\paragraph
\let\clearpage\relax
% Vista científica exacta materializada desde una rama condicional.
% Fuente: source/demostraciones_primarias/14_05b_extension_superviviente_caracter.tex; rama: HMTIncludeCharacterBlock.
% No modifica el contenido matemático de la rama de origen.
\chapter{Caractère positif de la signature et composition de registres}
\section{Caractère formel universel}

Soit \(\mathcal P=\RR^5\), de coordonnées
\[
 X=(X_A,X_C,X_4,X_{120},X_{270}).
\]
Pour \(v=(n_A,s_C,k,\nu_{120},\nu_{270})\in\ZZ^5\), nous définissons
\[
 \ell_{\rm form}(v;X)
 =n_AX_A+s_CX_C+kX_4
  +\nu_{120}X_{120}+\nu_{270}X_{270},
\]
\[
 \chi_{\rm form}(v;X)=\exp\!\bigl(\ell_{\rm form}(v;X)\bigr).
\]
Ainsi, \(\chi_{\rm form}\) prend ses valeurs dans le groupe multiplicatif
\(\operatorname{Fun}(\mathcal P,\RR_{>0})\), muni du produit ponctuel. Il n'a pas encore
sélectionné de coefficients globaux ni d'espèce.

\begin{theorem}[Caractère universel de la signature]
\label{thm:caracter-formal-ruta}
L'application
\[
 \chi_{\rm form}:(\ZZ^5,+)\longrightarrow
 \bigl(\operatorname{Fun}(\mathcal P,\RR_{>0}),\cdot\bigr)
\]
est un homomorphisme. La composition
\[
 \chi_{\rm route}^{\rm form}
 :=\chi_{\rm form}\circ L_{\rm tick}
\]
évalue formellement tout registre enrichi et n'introduit pas d'unité
physique.
\end{theorem}

\begin{proof}
Pour tout \(X\in\mathcal P\), l'exposant
\(\ell_{\rm form}(\,\cdot\,;X)\) est additif sur \(\ZZ^5\). Par conséquent,
\[
\chi_{\rm form}(v+w;X)
=\chi_{\rm form}(v;X)\chi_{\rm form}(w;X),
\]
et l'identité vaut comme égalité de fonctions positives.
\end{proof}

Le caractère formel précède la construction des coefficients globaux.
Cette section ne définit pas encore de caractère numérique de masse ni de
carte d'unités.

\section{Section interne minimale sur les 81 cellules}

L'atlas fini contient 81 cellules orientées. Leur projection sur la carte
de famille de chemin donne 56 valeurs. Notons cette projection
\[
 \mathcal R:\mathcal C_9\longrightarrow\mathcal R_{56}.
\]
L'histogramme exact de ses fibres est
\[
 41\times1,
 \qquad10\times2,
 \qquad2\times3,
 \qquad1\times4,
 \qquad2\times5,
\]
et leur somme est 81.

L'origine et l'ordre lexicographique de la carte étant fixés, soit
\(\kappa(x)\in\{0,1,2,3,4\}\) le rang de \(x\) au sein de sa fibre.

\begin{theorem}[Sélecteur interne minimal]
\label{thm:selector-interno-minimo}
L'application
\[
 x\longmapsto\bigl(\mathcal R(x),\kappa(x)\bigr)
\]
est injective sur les 81 cellules. Aucun alphabet de mémoire comportant moins de
cinq symboles ne peut relever \(\mathcal R\) injectivement sur l'ensemble de
l'atlas.
\end{theorem}

\begin{proof}
Au sein de chaque fibre, \(\kappa\) attribue des rangs distincts ; entre les fibres, la
première coordonnée est déjà distincte. Il existe une fibre de taille cinq, de sorte que
le principe des tiroirs impose à tout relèvement injectif de disposer d'au moins
cinq valeurs de mémoire. Le rang \(0,\ldots,4\) atteint cette borne.
\end{proof}

La canonicité est relative à l'origine et à l'ordre déclarés. \(\kappa\) n'est
ni une charge, ni une saveur, ni une masse. Les treize classes grossières de l'atlas produisent
des multisections finies ; seule la classe centrale possède le point fixe \((5,5)\)
sous l'inversion centrale. Les douze autres nécessitent une information orientée
pour sélectionner un point et n'admettent pas de sélecteur ponctuel équivariant à partir
de la seule classe grossière.

\subsection{Inversion de Möbius et caractères d'aller et de retour}

Sur le mot de signes dodécaphasique, nous définissons
\[
 C_M(\tau)_m=-\tau_{11-m}.
\]
Cette opération combine l'inversion de l'ordre de lecture avec le changement de
feuillet. Un caractère linéaire
\(\ell_a(\tau)=\sum_ma_m\tau_m\), avec
\(a_m=a_{11-m}\), satisfait
\[
 \ell_a(C_M\tau)=-\ell_a(\tau),
\]
tandis que toute forme quadratique symétrique compatible avec la réversion
satisfait
\[
 Q(C_M\tau)=Q(\tau).
\]
Ainsi, les coordonnées linéaires distinguent l'orientation d'aller et de retour et les
quadratiques conservent le contenu de corrélation. Cette parité explique
pourquoi \(s_C\) change de signe tandis que \(\nu_{270}\) peut rester
invariant. Dans l'atlas, les sept formes de mot ont pour multiplicités
\[
 54,2,2,2,7,7,7;
\]
c'est pourquoi l'antifeuillet se réalise comme correspondance entre fibres et non comme
une involution bijective obtenue en ignorant leurs multiplicités.

\section{Composition orientée de registres avant la projection}

Soient \(\Gamma\) et \(\Lambda\) deux histoires enrichies, \(B\) une
sous-histoire commune et \(g\) le transport qui identifie phase, feuillet,
orientation et position à l'interface. Au niveau du registre, on définit
\[
 \boxed{
 \Gamma\star_B\Lambda
 =\bigl(n_\Gamma+n_\Lambda-n_B,
        e_\Gamma+g e_\Lambda-e_B\bigr).}
\]
La définition présuppose aussi l'identification du bord, des
prédécesseurs torsionnels, de la couronne et de la mémoire antifeuillet.

\begin{proposition}[Obstruction à la descente pentadimensionnelle de la composition de registres]
\label{prop:no-descenso-empalme-cinco}
En général, il n'existe pas d'opération binaire sur \(\ZZ^5\) qui récupère
\(L_{\rm tick}(\Gamma\star_B\Lambda)\) à partir de
\(L_{\rm tick}(\Gamma)\) et \(L_{\rm tick}(\Lambda)\) seulement.
\end{proposition}

\begin{proof}
L'opération compose les douze événements avant d'appliquer
\(\operatorname{sgn}\) et les autocorrélations. Deux histoires peuvent avoir
le même quintuplet et des valeurs distinctes de \(M_5\) ou \(A_6\). En les composant
avec une troisième histoire, une somme
\(e_a+e_{a+4}+e_{a+8}\) peut franchir zéro dans l'une et non dans l'autre, modifiant
\(\nu_{120}^{\rm ev}\). Les losanges finis fournissent des témoins explicites.
Par conséquent, la projection sur cinq entiers n'est pas
une congruence pour \(\star_B\).
\end{proof}

La proposition formalise une règle opératoire essentielle : les extensions se
composent ; les signatures se lisent ensuite. La particule composée ne
s'obtient pas non plus en additionnant aveuglément deux lignes d'un tableau, mais en composant leurs mémoires
compatibles et en évaluant le résultat.

\section{Dépendance dodécaphasique du caractère de persistance}

Les moteurs réduits de période 54 imposent
\[
 e_{m+6}=e_m.
\]
Il en découle les congruences
\[
 n_A\equiv0\pmod2,
 \qquad
 s_C\equiv0\pmod2,
 \qquad
 \nu_{270}\equiv0\pmod4.
\]
Les signatures de la reconstruction scalaire abrégée violent au moins une de
ces congruences. Leur généalogie exige donc le registre enrichi et
ne procède pas de ce moteur réduit avec les règles déclarées.

Ce résultat n'est pas une obstruction à la dynamique TPK enrichie. Il démontre
quelle mémoire elle nécessite : une dodécaphase véritable, avec éventuellement
\(A_6\ne0\), une composition orientée et des prédécesseurs torsionnels conservés. L'ablation
évite d'attribuer à TPK une limitation qui appartient à une réduction
périodique concrète.

\section{Factorisation par le quotient de registre et réalisation physique}

Dans le cadre de cette section, on déclare la relation
\[
\gamma\sim_{\mathcal L}\gamma'
\quad\Longleftrightarrow\quad
\text{\(\gamma,\gamma'\) présentent la même extension et }
\mathcal L(\gamma)=\mathcal L(\gamma').
\]
Nous notons
\(\mathcal E_{108}^{\rm surv}:=
\Gamma_{108}^{\rm enr}/\!\sim_{\mathcal L}\)
le quotient relatif à cette carte. Soit
\[
 \pi_{\rm ext}:\Gamma_{108}^{\rm enr}\longrightarrow
 \mathcal E_{108}^{\rm surv}
\]
la projection canonique. Par construction, l'extracteur est constant sur ses
fibres. L'affirmation plus forte selon laquelle toute carte admissible d'une
même extension produit automatiquement le même registre exigerait de démontrer
la compatibilité entre cartes et n'est pas présupposée ici.

\begin{theorem}[Factorisation relative par le quotient de registre]
\label{thm:descenso-caracter-extension}
Il existe une unique application
\[
 \overline L:\mathcal E_{108}^{\rm surv}\longrightarrow\ZZ^5
\]
tel que
\[
 \boxed{L_{\rm tick}=\overline L\circ\pi_{\rm ext}.}
\]
Par conséquent, \(\chi_{\rm form}\circ\overline L\) est une fonction de la classe dans
le quotient relatif et non du représentant de chemin choisi pour la lire.
\end{theorem}

\begin{proof}
Si \(\pi_{\rm ext}(\gamma)=\pi_{\rm ext}(\gamma')\), les deux présentations
possèdent le même registre par la relation déclarée et, d'après les formules
précédentes, le même quintuplet. Nous définissons
\(\overline L([\gamma])=L_{\rm tick}(\gamma)\). La constance prouve que la
définition est indépendante du représentant ; la surjectivité de l'application
au quotient prouve l'unicité.
\end{proof}

La construction interne atteint ainsi
\[
 \Omega_{\HMT}^{\rm surv}
 \longrightarrow\mathcal E_{108}^{\rm surv}
 \xrightarrow{\ \overline L\ }\ZZ^5
 \xrightarrow{\ \chi_{\rm form}\ }
 \operatorname{Fun}(\mathcal P,\RR_{>0}).
\]
Le contenu du théorème est positif et complet dans ce domaine : une
extension survivante détermine sa classe de registre, la classe détermine
une unique signature et la signature détermine le caractère formel. La spécialisation
globale et la hiérarchie complète des tours sont construites dans les chapitres
suivants ; les
comparaisons avec les noms et les unités physiques sont présentées ensuite, sans
modifier cette chaîne.

\endgroup
\begingroup
\let\chapter\section
\let\section\subsection
\let\subsection\subsubsection
\let\subsubsection\paragraph
\let\clearpage\relax
\chapter{Algèbre déployée et factorisation exacte du caractère}
\label{chap:factorizacion-split}
\index{algèbre déployée}\index{caractère!factorisation déployée}
\index{cône nul}\index{canaux $q_\pm$}

La signature formelle d'une extension et les canaux angulaires proviennent de branches déjà construites. Dans cette section, ils convergent vers une factorisation déterminantielle qui conserve séparément l'échelle radiale et l'orientation relative. Le résultat appartient au niveau algébrique : il n'identifie à lui seul ni les idéaux nuls à des photons physiques ni l'intérieur positif à une particule matérielle.

\section{Coordonnées angulaires conjuguées et réseau spectral bidirectionnel}
\label{sec:reticula-angular-split}

La chaîne causale précédente fixe
\[
 A_{\rm rad}=\frac{\pi A}{180},\qquad
 \Cstarrad=\frac{\pi\Cstar}{180},
\]
et, seulement après l'obtention de la coordonnée directe \(A\) et de l'unique coordonnée conjuguée \(C^*\),
\[
 q_+=\exp[-(A_{\rm rad}+\Cstarrad)],
 \qquad
 q_-=\exp[-(A_{\rm rad}-\Cstarrad)].
\]
Pour une signature angulaire \((n_A,s_C)\), introduisons les deux nombres d'enroulement
\[
 W_+=6n_A+s_C,
 \qquad
 W_-=6n_A-s_C.
\]

\begin{theorem}[Réseau angulaire d'indice douze]
\label{thm:reticula-two-way}
L'application
\[
 \mathcal W:\ZZ^2\to\ZZ^2,
 \qquad
 \binom{n_A}{s_C}\longmapsto
 \binom{W_+}{W_-}
 =\begin{pmatrix}6&1\\6&-1\end{pmatrix}
 \binom{n_A}{s_C}
\]
est injective et son image est d'indice douze. Ses facteurs invariants sont \((1,12)\), et
\[
 \operatorname{im}\mathcal W
 =\{(u,v)\in\ZZ^2:u+v\equiv0\pmod{12}\}.
\]
De plus,
\[
 \boxed{
 \exp\!\left(n_AA_{\rm rad}+\frac{s_C}{6}\Cstarrad\right)
 =\left(q_+^{-W_+}q_-^{-W_-}\right)^{1/12}.}
\]
\end{theorem}

\begin{proof}
La matrice a pour déterminant \(-12\). Le plus grand commun diviseur de ses coefficients étant égal à un, sa forme normale de Smith est \(\operatorname{diag}(1,12)\). Si \((u,v)=(6n+s,6n-s)\), alors \(u+v=12n\). Réciproquement, si \(u+v\) est un multiple de douze, alors
\[
 n=\frac{u+v}{12},
 \qquad
 s=\frac{u-v}{2}
\]
sont entiers et permettent de retrouver \((u,v)\). Enfin,
\[
 \frac{W_+(A_{\rm rad}+\Cstarrad)
       +W_-(A_{\rm rad}-\Cstarrad)}{12}
 =n_AA_{\rm rad}+\frac{s_C}{6}\Cstarrad,
\]
ce qui démontre l'identité par exponentiation.
\end{proof}

L'involution \(s_C\mapsto-s_C\) échange \(W_+\leftrightarrow W_-\) et \(q_+\leftrightarrow q_-\). La racine douzième correspond à l'indice entier de l'image de \(\mathcal W\), et non à un arrondi. La coordonnée \(C^*\) agit comme composante antisymétrique du réseau ; ce n'est ni une correction décimale d'une masse ni une sortie de \(\Gamma_{6\to8}\).

\section{Algèbre réelle déployée}
\label{sec:algebra-real-partida}

\begin{definition}[Algèbre réelle déployée]
Soit
\[
 \mathcal A_{\rm sp}
 :=\RR[\mathsf R]/(\mathsf R^2-1).
\]
La conjugaison et la norme sont définies respectivement par
\[
 \overline{a+b\mathsf R}=a-b\mathsf R,
 \qquad
 N_{\rm sp}(z)=z\overline z=a^2-b^2.
\]
\end{definition}

Les idempotents conjugués
\[
 P_\pm=\frac{1\pm\mathsf R}{2}
\]
satisfont
\[
 P_\pm^2=P_\pm,\qquad
 P_+P_-=0,\qquad
 P_++P_-=1,\qquad
 \overline{P_+}=P_-.
\]

\begin{proposition}[Décomposition spectrale déployée]
\label{prop:descomposicion-espectral-partida} Tout \(z\in\mathcal A_{\rm sp}\) admet une décomposition unique
\[
 z=r_+P_++r_-P_-,
 \qquad r_\pm\in\RR,
\]
et sa norme est
\[
 \boxed{N_{\rm sp}(z)=r_+r_-.}
\]
Chaque idéal principal \(\mathcal A_{\rm sp}P_\pm=\RR P_\pm\) est une droite nulle ; en particulier,
\[
 N_{\rm sp}(\lambda P_\pm)=0.
\]
\end{proposition}

\begin{proof}
Si \(z=a+b\mathsf R\), le choix \(r_+=a+b\), \(r_-=a-b\) fournit la décomposition, dont l'unicité découle de l'indépendance de \(P_+\) et \(P_-\). La conjugaison échange les deux coefficients : \(\overline z=r_-P_++r_+P_-\). L'orthogonalité des idempotents donne
\[
 z\overline z=(r_+r_-)(P_++P_-)=r_+r_-.
\]
L'énoncé relatif à chaque idéal s'obtient en annulant l'un des deux coefficients.
\end{proof}

\begin{definition}[Cône causal déployé]
\label{def:cono-causal-partido} Le cône positif, sa frontière nulle et son intérieur sont
\[
 \mathcal C_{\rm sp}^{+}
 =\{r_+P_++r_-P_-:r_+,r_-\geq0\},
\]
\[
 \partial_{\rm null}\mathcal C_{\rm sp}^{+}
 =\RR_{\geq0}P_+\cup\RR_{\geq0}P_-,
\qquad
 \operatorname{int}\mathcal C_{\rm sp}^{+}
 =\{r_+P_++r_-P_-:r_+r_->0\}.
\]
\end{definition}

La frontière possède une seule direction active et une norme nulle. L'intérieur exige deux directions simultanées et une norme positive. Cette stratification est exacte ; sa réalisation physique exige des applications qui préservent la norme, le transport et la dynamique concernés.

\section{Transporteur relatif et conditions de réalisation représentationnelle}
\label{sec:transportador-relativo-split-md}

Le bloc historique \(B_0=7I+2\mathsf R\) a pour spectre \(\{9,5\}\), et le bloc angulaire \(B_1=AI+\Cstar\mathsf R\) a pour spectre \(\{A+\Cstar,A-\Cstar\}\). Le \cref{thm:no-go-entrelazador-bloques-split} démontre qu'avec les valeurs internes en vigueur, tout entrelaceur linéaire supposé \(\Phi B_0=B_1\Phi\) est nul. L'objet actif est le transporteur multiplicatif relatif déjà construit, et non une équivalence entre représentations de spectres disjoints.

\section{Factorisation déployée du caractère angulaire}
\label{sec:factorizacion-split-angular}

Soit
\[
 \sigma=\frac{s_C}{6},\qquad
 w_\pm=\frac{n_A\pm\sigma}{2}=\frac{W_\pm}{12},
 \qquad
 r_\pm=q_\pm^{-w_\pm},
\]
et définissons
\[
 Z_v=r_+P_++r_-P_-.
\]

\begin{theorem}[Factorisation déployée du caractère angulaire]
\label{thm:split-character} Pour toute signature angulaire \(v=(n_A,s_C)\in\ZZ^2\),
\[
 \boxed{
 N_{\rm sp}(Z_v)
 =\det Z_v
 =\exp\!\left(
 n_AA_{\rm rad}+\frac{s_C}{6}\Cstarrad
 \right).}
\]
\end{theorem}

\begin{proof}
Par définition des canaux,
\[
 \log r_+
 =\frac{n_A+\sigma}{2}(A_{\rm rad}+\Cstarrad),
 \qquad
 \log r_-
 =\frac{n_A-\sigma}{2}(A_{\rm rad}-\Cstarrad).
\]
En sommant,
\[
 \log(r_+r_-)
 =n_AA_{\rm rad}+\sigma\Cstarrad.
\]
La \cref{prop:descomposicion-espectral-partida} donne \(N_{\rm sp}(Z_v)=r_+r_-\), et la représentation diagonale dans la base \((P_+,P_-)\) donne la même expression pour le déterminant.
\end{proof}

La composante que le déterminant ne détecte pas conserve l'orientation relative :
\[
 \frac12\log\frac{r_+}{r_-}
 =\frac12\left(
 n_A\Cstarrad+\frac{s_C}{6}A_{\rm rad}
 \right),
\]
\[
 Z_v=
 \exp\!\left[
 \frac12\left(n_AA_{\rm rad}+\frac{s_C}{6}\Cstarrad\right)
 \right]
 \exp\!\left[
 \frac12\left(n_A\Cstarrad+\frac{s_C}{6}A_{\rm rad}\right)
 \mathsf R
 \right].
\]
Le premier facteur fixe l'échelle radiale ; le second enregistre l'orientation relative. Ce sont deux lectures coordonnées d'un même élément.

\section{Incorporation centrale des tours}
\label{sec:factorizacion-split-torres}

Une fois \(\Dfour\), \(s_{120}\) et \(\zCorr=\zeta_{270}^{\rm corr}\) construits globalement, soit
\[
 t(v)=
 k_{\Dfour}\Dfour
 +\nu_{120}s_{120}
 +\nu_{270}\zCorr,
 \qquad
 \widetilde Z_v=e^{t(v)/2}Z_v.
\]

\begin{corollary}[Caractère complet comme norme déployée]
\label{cor:split-character-torres} Les blocs globaux précédents étant fixés,
\[
 \boxed{
 N_{\rm sp}(\widetilde Z_v)
 =\exp\!\left(
 n_AA_{\rm rad}+\frac{s_C}{6}\Cstarrad
 +k_{\Dfour}\Dfour+\nu_{120}s_{120}
 +\nu_{270}\zCorr
 \right).}
\]
\end{corollary}

\begin{proof}
Dans la représentation bidimensionnelle, la multiplication par le scalaire central \(e^{t(v)/2}\) multiplie le déterminant par \(e^{t(v)}\). On applique le \cref{thm:split-character}.
\end{proof}

Le corollaire construit la réalisation déployée déclarée du caractère actif ; il ne classe pas toutes les réalisations possibles et ne transforme pas les tours en preuves indépendantes. Ses hypothèses sont les canaux et les blocs globaux déjà fixés, et son contrôle négatif est immédiat : omettre le facteur central \(1/2\) double l'exposant de tour ; activer simultanément les deux idempotents avec des coefficients positifs fait sortir de la frontière nulle.

\endgroup
\begingroup
\let\chapter\section
\let\section\subsection
\let\subsection\subsubsection
\let\subsubsection\paragraph
\let\clearpage\relax
\chapter{Valeurs, complétude et blocs de la hiérarchie spectrale}
\label{chap:torres}

L'opérateur à deux couches, ses moments et le semi-groupe \(\mathfrak S=\langle90,120\rangle\) sont définis une seule fois dans le \cref{sec:jerarquia-espectral-infinita-rev7}. Ce chapitre expose des valeurs, démontre la complétude reconstructive et sépare les blocs globaux. Les mêmes entiers \(90,120\) apparaissent également comme cardinalités dans l'incidence hexade--octade, mais cette coïncidence ne transporte pas l'incidence \(\iota_{6,8}\) dans la carte angulaire et ne transforme pas ses configurations en moments de masse.

\section{Premières évaluations de la hiérarchie}

Les indices
\[
 90,120,180,210,240,270,360,420
\]
constituent une sélection de premières hauteurs de \(\mathfrak S\), puisque
\[
 180=2\cdot90,\quad210=90+120,\quad240=2\cdot120,
\]
\[
 270=3\cdot90,\quad360=3\cdot120,
 \quad420=2(90+120).
\]
Les moments \(s_n=q_-^n-q_+^n\) et l'interface génératrice \(\mathcal K_{\mathrm{ang}}^{90,120}=(s_{90},s_{120})\) sont utilisés ici selon la convention déjà fixée ; ils ne sont pas redéfinis. L'ensemble des valeurs est déterminé par \(A\), \(\Cstar\) et la récurrence quadratique du lecteur.

\begin{table}[htbp]
\centering
\caption{Moments orientés de la hiérarchie globale.}
\label{tab:torres}
\begin{tabular}{@{}rr@{}}
\toprule
\(n\)&\(s_n\)\\
\midrule
90  & \(2.9516943928982796\times10^{-4}\)\\
120 & \(1.9683511250294992\times10^{-5}\)\\
180 & \(8.7345871869508335\times10^{-8}\)\\
210 & \(5.8181313057291163\times10^{-9}\)\\
240 & \(3.8754674969305336\times10^{-10}\)\\
270 & \(2.5814553607509555\times10^{-11}\)\\
360 & \(7.6293257871406076\times10^{-15}\)\\
420 & \(3.3850663628348497\times10^{-17}\)\\
\bottomrule
\end{tabular}
\end{table}

\begin{proposition}[Caractère global des moments]
Les quantités \(s_n\) de la \cref{tab:torres} sont des invariants de la carte angulaire fixée. Une espèce fournit des coordonnées entières qui pondèrent ces invariants, mais ne modifie pas \(q_\pm\) et n'introduit pas de tour propre.
\end{proposition}
\begin{proof}
Chaque \(s_n\) est une fonction de \(q_-\) et de \(q_+\). Le théorème d'inversion des canaux démontre que l'application conjointe \((A,\Cstar)\mapsto(q_+,q_-)\) est injective. Aucune étiquette d'espèce n'apparaît dans ces définitions.
\end{proof}

\begin{theorem}[Complétude reconstructive de la hiérarchie]
\label{thm:completitud-reconstructiva-torres} Soit \(n_1<n_2<\cdots\) l'énumération croissante sans répétitions de \(\mathfrak S\). Pour
\[
 a_k:=s_{n_k}=q_-^{n_k}-q_+^{n_k}>0,
\]
tout rapport positif \(\rho\in\RR_{>0}\) admet un développement absolument convergent dans la hiérarchie :
\[
 \boxed{\rho=
 \exp\!\left(\sum_{k\geq1}\nu_k a_k\right)}
 \qquad\text{pour une certaine suite }(\nu_k)_{k\geq1}\subset\ZZ.
\]
\end{theorem}

\begin{proof}
Puisque \(0<q_+<q_-<1\), on a
\[
 0<a_k<q_-^{n_k}.
\]
Les \(n_k\) sont des multiples de \(30\) supérieurs ou égaux à \(90\) ; donc,
\[
 \sum_{k\geq1}a_k
 \leq\sum_{m\geq3}q_-^{30m}
 =\frac{q_-^{90}}{1-q_-^{30}}<\infty,
 \qquad a_k\longrightarrow0.
\]
Fixons \(r_0=\log\rho\) et une règle de départage pour l'entier le plus proche. Définissons récursivement
\[
 \nu_k=\operatorname{nint}\!\left(\frac{r_{k-1}}{a_k}\right),
 \qquad
 r_k=r_{k-1}-\nu_k a_k.
\]
La propriété de l'entier le plus proche donne
\[
 |r_k|\leq\frac{a_k}{2}\longrightarrow0.
\]
L'identité télescopique
\[
 r_0=\sum_{j=1}^{N}\nu_j a_j+r_N
\]
implique \(r_0=\sum_{j\geq1}\nu_j a_j\). La convergence est absolue, puisque
\[
 \begin{aligned}
 \sum_{k\geq1}|\nu_k|a_k
 &=\sum_{k\geq1}|r_{k-1}-r_k|\\
 &\leq |r_0|+\frac{a_1}{2}
 +\frac12\sum_{k\geq2}(a_{k-1}+a_k)<\infty.
 \end{aligned}
\]
L'exponentiation donne l'énoncé.
\end{proof}

\begin{remark}[Portée et contrôle négatif]
Le théorème établit une surjectivité reconstructive du langage de tours sur \(\RR_{>0}\) : il démontre la résolution du codomaine, et non la sélection physique. L'algorithme reçoit \(r_0=\log\rho\) ; si \(\rho\) est une grandeur cible, ses coefficients constituent donc une reconstruction dépendant de cette cible. La provenance prospective exige une autre flèche,
\[
 \text{espèce}\dashrightarrow\text{extension}
 \longrightarrow\text{route}\longrightarrow\text{registre}
 \longrightarrow\text{signature},
\]
et ne se déduit pas de ce développement. La conclusion serait invalidée si les \(a_k\) ne tendaient pas vers zéro, si la borne du résidu cessait d'être satisfaite ou si la série des erreurs n'était pas sommable.
\end{remark}

\section{Blocs globaux de l'action de masse}

L'action de masse utilisera le triplet
\[
 \Dfour,\qquad s_{120},\qquad \zCorr=135\Dfour^2.
\]
Leurs valeurs sont calculées une seule fois. Les coordonnées \(k,\nu_{120},\nu_{270}\) d'une signature précisent le nombre d'interventions de chaque bloc dans l'exposant ; elles n'en modifient pas la valeur.

Il convient de distinguer deux objets qui partagent l'indice \(270\) :
\[
 \boxed{\zCorr=135\Dfour^2
 =1.66922699663643\times10^{-6}}
\]
et
\[
 \boxed{\sSpec=q_-^{270}-q_+^{270}
 =2.58145536075096\times10^{-11}.}
\]
Le premier est un invariant quadratique construit à partir de \(\Dfour\) ; le second est un moment spectral du semi-groupe. La coïncidence de l'indice n'autorise pas leur identification.

\section{Combinaison d'ordre supérieur}

La combinaison globale
\[
 \delta\Dfour^{\rm tower}
 =\frac{s_{180}}{24}-\frac{s_{240}}{54}
 -24s_{360}+\frac76s_{420}
\]
a pour valeur
\[
 \delta\Dfour^{\rm tower}
 =3.632051471908759123889070967\times10^{-9}.
\]
La coordonnée de comparaison utilisée pour sélectionner les coefficients est définie comme la donnée scalaire
\[
 \delta\Dfour
 :=3.632051471908972784661641720\times10^{-9}.
\]
Par conséquent, le résidu de la combinaison par rapport à cette coordonnée est
\[
 \delta\Dfour^{\rm tower}-\delta\Dfour
 =-2.1366077257075\times10^{-22}.
\]
La petite quantité est donc une différence entre deux objets, et non la valeur de la combinaison. Tous les termes appartiennent à la même hiérarchie et leurs coefficients sont globaux. La coordonnée de comparaison étant intervenue dans la sélection des coefficients, cette égalité est reconstructive. La combinaison n'est pas ajoutée une nouvelle fois à la formule électronique qui emploie déjà \(\Dfour\), car cette substitution compterait deux fois la même correction.

\endgroup
\begingroup
\let\chapter\section
\let\section\subsection
\let\subsection\subsubsection
\let\subsubsection\paragraph
\let\clearpage\relax
\section{Opérateur à deux couches et hiérarchie infinie}
\label{sec:jerarquia-espectral-infinita-rev7}
\index{tours!hiérarchie infinie}\index{semi-groupe spectral}

Soient
\[
 x=A_{\rm rad},\qquad y=C^*_{\rm rad},\qquad
 R_{\rm ch}^2=I,\qquad
 P_\pm^{\rm ch}=\frac{I\pm R_{\rm ch}}2.
\]
Le lecteur des deux feuillets est réuni dans l'opérateur
\[
 T_{\rm ang}=\exp(-xI-yR_{\rm ch})
 =q_+P_+^{\rm ch}+q_-P_-^{\rm ch},
\]
avec
\[
 q_-=e^{-x+y},\qquad q_+=e^{-x-y}.
\]
Ses moments spectraux pair et impair sont
\[
 c_n=\operatorname{Tr}(T_{\rm ang}^n)=q_-^n+q_+^n,\qquad
 s_n=-\operatorname{Tr}(R_{\rm ch}T_{\rm ang}^n)=q_-^n-q_+^n.
\]
L'involution \(C^*\mapsto-C^*\) conserve \(c_n\) et inverse \(s_n\). De manière équivalente,
\[
 c_n=2e^{-nx}\cosh(ny),\qquad
 s_n=2e^{-nx}\sinh(ny).
\]
Il s'agit de la convention fixée dans le chapitre consacré aux canaux. Pour déclarer sans homonymie la présentation alternative des sources relatives aux tours, posons
\[
 R_{\rm tow}:=-R_{\rm ch},\qquad
 P_\pm^{\rm tow}=P_\mp^{\rm ch},
\]
\[
 T_{\rm tow}:=e^{-xI+yR_{\rm tow}}=T_{\rm ang},\qquad
 \operatorname{Tr}(R_{\rm tow}T_{\rm tow}^n)
 =-\operatorname{Tr}(R_{\rm ch}T_{\rm ang}^n)=s_n.
\]
Les deux écritures décrivent le même lecteur scalaire et ne constituent pas une seconde orientation des tours.

\begin{theorem}[Addition, norme et récurrence]
\label{thm:recurrencia-bicapa-rev7}
Pour \(m,n\ge0\),
\[
 c_{m+n}=\frac{c_mc_n+s_ms_n}{2},\qquad
 s_{m+n}=\frac{s_mc_n+c_ms_n}{2},
\]
\[
 c_n^2-s_n^2=4e^{-2nx}.
\]
Si \(p=q_-q_+=e^{-2x}\), les deux suites satisfont
\[
 X_{n+2}=c_1X_{n+1}-pX_n.
\]
De plus,
\[
 \partial_Ac_n=-\frac{n\pi}{180}c_n,
 \qquad
 \partial_As_n=-\frac{n\pi}{180}s_n,
\]
\[
 \partial_{C^*}c_n=\frac{n\pi}{180}s_n,
 \qquad
 \partial_{C^*}s_n=\frac{n\pi}{180}c_n.
\]
\end{theorem}

\begin{proof}
Les identités d'addition résultent du développement de \((q_-^m\pm q_+^m)(q_-^n\pm q_+^n)\). La norme est
\[
 (q_-^n+q_+^n)^2-(q_-^n-q_+^n)^2
 =4(q_-q_+)^n.
\]
La dernière égalité est la récurrence déterminée par les racines caractéristiques \(q_-\) et \(q_+\). Les équations différentielles s'obtiennent en dérivant les expressions hyperboliques ; le facteur \(\pi/180\) appartient à la carte angulaire, et non à une sélection d'espèce.
\end{proof}

\begin{theorem}[Semi-groupe complet des hauteurs]
\label{thm:semigrupo-alturas-rev7} Les hauteurs de la hiérarchie transversale forment
\[
 \mathfrak S
 =\langle90,120\rangle
 =\{90a+120b:a,b\in\mathbb Z_{\ge0},\ a+b>0\}.
\]
De manière équivalente,
\[
 \boxed{\mathfrak S=\{90,120\}\cup\{30r:r\ge6\}.}
\]
\end{theorem}

\begin{proof}
En divisant par \(30\), on obtient le semi-groupe \(\langle3,4\rangle\). Ses éléments positifs sont \(\{3,4\}\cup\{r:r\ge6\}\) : \(5\) est l'unique lacune positive postérieure aux générateurs, et tout \(r\ge6\) s'obtient par récurrence en ajoutant \(3\) ou \(4\). En multipliant de nouveau par \(30\), on obtient la formule.
\end{proof}

La présentation conserve également la généalogie d'une hauteur :
\[
 \mathfrak S\cong
 \bigl(\mathbb Z_{\ge0}^2\setminus\{(0,0)\}\bigr)
 \big/\bigl((a+4,b)\sim(a,b+3)\bigr),
\]
puisque \(4\cdot90=3\cdot120=360\). La hauteur ne détermine pas à elle seule la composition qui l'a produite ; les généalogies \(4\cdot90\) et \(3\cdot120\) restent distinctes dans le registre chronologique enrichi, bien qu'elles donnent le même indice.

Pour expliciter la récurrence échantillonnée au pas commun de trente, définissons
\[
 a_{30}:=q_-^{30},\qquad b_{30}:=q_+^{30},\qquad
 u_r:=s_{30r}=a_{30}^r-b_{30}^r.
\]
Alors
\[
 \boxed{u_{r+2}=(a_{30}+b_{30})u_{r+1}-a_{30}b_{30}u_r.}
\]
Le terme auxiliaire \(u_5=s_{150}\) peut intervenir dans la propagation de la récurrence, mais \(150\notin\mathfrak S\) ; une identité de calcul ne crée pas de nouvelle hauteur admissible.

En particulier,
\[
90,120,180,210,240,270,300,330,360,390,420,450,480,\ldots
\]
sont les premières hauteurs. Les huit hauteurs mises en évidence dans les exposés antérieurs sont des premières lectures de \(\mathfrak S\), et non sa définition ou sa limite.

À titre de guide de lecture, sa généalogie initiale peut se lire sans être assimilée au domaine complet :
\[
\begin{array}{c@{\quad}c@{\quad}l}
90&3\cdot30&\text{générateur minimal}\\
120&4\cdot30&\text{générateur minimal}\\
180&2\cdot90&\text{premier mode pair}\\
210&90+120&\text{mode composé}\\
240&2\cdot120&\text{second mode pair}\\
270&3\cdot90&\text{répétition de la branche 90}\\
360&3\cdot120&\text{répétition de la branche 120}\\
420&2(90+120)&\text{retour composé}
\end{array}
\]
La dernière colonne consigne la provenance structurelle ; elle n'introduit pas de paramètres et n'affirme pas que ces huit lignes épuisent la hiérarchie. Les mêmes entiers \(90,120\) sont des cardinalités d'incidence dans \(\mathcal F_6\) et \(\mathcal F_8\), mais cette coïncidence n'identifie pas l'opérateur à deux couches à l'inclusion \(\iota_{6,8}\).

L'interface des deux générateurs se réduit au couple
\[
 \mathcal K_{\rm ang}^{90,120}(q_+,q_-):=(s_{90},s_{120}).
\]
Elle ne doit être confondue ni avec les résolvantes de deux sous-queues ni avec la fonctionnelle ultérieure \(K_{6\to8}=12\Delta S_{90}-\Delta S_{120}\), dont la définition et la démonstration se trouvent dans le \cref{mdg:chap:barbero}. Les sous-queues y appartiennent à la complétion analytique de la hiérarchie ; elles ne sont ni de nouveaux générateurs ni des coefficients d'espèce.

\begin{definition}[Module convergent de raffinements spectraux]
\label{def:modulo-refinamientos-torre-rev7} Le domaine intégral du relèvement global est
\[
 \mathcal E_{\mathfrak S}
 :=\left\{\eta\in\mathbb Z^{\mathfrak S}:
 \sum_{h\in\mathfrak S}|\eta_hs_h|<\infty\right\},
 \qquad s_h=q_-^h-q_+^h.
\]
Chaque \(\eta\in\mathcal E_{\mathfrak S}\) détermine le caractère spectral
\[
 \mathcal R_\eta
 :=\exp\!\left(\sum_{h\in\mathfrak S}\eta_hs_h\right).
\]
\end{definition}

La loi des masses ne réside pas dans cette définition : sa signature centrale \((n_A,s_C,k_{\Delta_4},b_{120},b_\zeta)\in\mathbb Z^5\) est construite auparavant à partir du registre. Le relèvement par \(\eta\) agit ensuite et conserve la provenance du terme de base \(b_{120}s_{120}\). En particulier, \(N_{120}=b_{120}+\eta_{120}\) est le coefficient total évalué, et non une nouvelle coordonnée de signature ou une décomposition inverse unique.

La complétude du semi-groupe ne sélectionne pas à elle seule une espèce. Lorsque \(\eta\) est obtenu à partir d'un résidu métrologique fourni, il s'agit d'une reconstruction calibrée exacte ; la portée prospective exige que la route, la retenue et l'enroulement produisent \(\eta\) avant l'examen de ce résidu. La récurrence précédente engendre tous les moments une fois le raffinement fixé, sans introduire une loi distincte à chaque hauteur.

\endgroup
\begingroup
\let\chapter\section
\let\section\subsection
\let\subsection\subsubsection
\let\subsubsection\paragraph
\let\clearpage\relax
\section{Paire de lecteurs généalogiques sur l'atlas complet}
\label{sec:operadores-masa-two-way-rev11}
\index{masse!opérateurs bidirectionnels}
\index{atlas!lecteurs chronologique et dodécaphasique}

Le parcours enrichi possède deux représentations internes complémentaires,
\[
 \Gamma\longmapsto
 \bigl(\gamma_{108}(\Gamma),\tau_{12}(\Gamma)\bigr)
 \in\mathbb F_3^{108}\times\{-1,0,+1\}^{12}.
\]
La première conserve l'histoire tritique complète ; la seconde conserve
l'enregistrement dodécaphasique signé. Chaque projection induit un lecteur naturel, et les deux
sont évalués avant l'introduction des masses, des étiquettes PDG ou des résidus.

\HMTClaim{MD-R-MASS-TWOWAY-READER-PAIR}
\subsection{Paire de lecteurs généalogiques \(\mathbf K(\Gamma)=(K_D(\Gamma),K_\Omega(\Gamma))\) : domaine, codomaine et évaluations scalaires}

Avec la notation des
\cref{eq:KD-rev11,eq:kappaD-rev11,eq:kappa-omega-rev11}, nous définissons
\begin{equation}
 \mathbf K(\Gamma)
 :=\bigl(K_D(\Gamma),K_\Omega(\Gamma)\bigr)
 \in\left(\mathbb Z^2\times\tfrac12\mathbb Z\right)
 \times\mathbb Z^3,
 \label{eq:lector-conjunto-masas-rev11}
\end{equation}
où
\[
 K_D=\left(D_{27}+D_{36},D_1,\frac{D_4-D_3}{2}\right),
 \qquad
 K_\Omega=(\Omega_{90\mid120},54,P_6).
\]
Les évaluations scalaires sont
\begin{align}
 \kappa_D(\Gamma)
 &=D_{27}+D_{36}-\alpha_{\rm HMT}D_1
 +\frac{\Delta_4}{2}(D_4-D_3),
 \label{eq:kappaD-general-rev11}\\
\kappa_\Omega(\Gamma)
&=\Omega_{90\mid120}(\tau_\Gamma)
-54\alpha_{\rm HMT}+P_6(\tau_\Gamma)\Delta_4.
\label{eq:kappaOmega-general-rev11}
\end{align}
La deuxième coordonnée exige une distinction causale. Dans APP, le saut
spectral primordial est \(9-5=4\) ; le coefficient local \(2\), la compression
\(6=51-45\) et le déploiement bidimensionnel \(54=9\cdot6=459-405\) sont
des invariants ultérieurs et typés, conformément au
\cref{thm:app-jerarquia-cuatro-dos-seis-cincuentaycuatro}. Dans le lecteur
chronologique des déplacements, en revanche, la coordonnée électronique est
\(B_D=D_1(\Gamma_e)=54\) : elle compte les frontières immédiates et coïncide avec la
périodicité cinématique d'ordre \(54\) de l'endomorphisme CT108. Le
\(54\) fixe de
\(K_\Omega\) est la coordonnée commune transférée lorsque l'on rend compatibles les deux
lecteurs dans la réduction électronique. Son égalité numérique avec le défaut
global d'APP constitue un contrôle croisé ultérieur ; elle ne fait pas du
nombre \(54\) la différence primordiale entre somme et produit.

L'égalité \(4\cdot90=3\cdot120=360\) explique la généalogie commune des
deux lecteurs. \(K_D\) compare les déplacements du mot tritique
\(\gamma_{108}\) ;
\(K_\Omega\) compare les énergies de fenêtres dans le mot de douze phases. La
coïncidence d'origine n'identifie pas leurs quotients d'information.
Dans le domaine certifié \(\Gamma_{324}^{\rm or}\), la périodicité de
demi-tour rend pairs tous les \(D_k\), de sorte que
\((D_4-D_3)/2\in\mathbb Z\) et \(K_D\in\mathbb Z^3\). Hors de ce domaine,
la troisième composante conserve le type \(\tfrac12\mathbb Z\).

\subsection{Naturalité des opérateurs bidirectionnels et préservation des types}

Le lecteur \(D_k\) est invariant par translation temporelle, renversement et toute
permutation bijective de l'alphabet tritique. Dans les parcours déclarés,
l'inversion simultanée des deux orientations agit par une translation de
27 pas discrets et conserve \(K_D\). Le lecteur \(K_\Omega\) est invariant par
rotation et renversement dodécaphasiques, ainsi que par
\(C_M(\tau)=-\operatorname{rev}(\tau)\). Sa forme quadratique
\[
 \Omega(\tau)=\langle\tau,
 (4W_3^*W_3-3W_4^*W_4)\tau\rangle
\]
possède une polarisation bilinéaire entière et satisfait l'identité de 2-cocycle.

Les compteurs globaux de rotation, de changement de feuillet, de répétition de phase, de changement de
préfixe et de retour sont constants sur les 324 parcours. Par conséquent, toute
fonctionnelle qui se factorise exclusivement par ces six compteurs est constante
et se simplifie dans les rapports. Les lecteurs de
\eqref{eq:kappaD-general-rev11} et de
\eqref{eq:kappaOmega-general-rev11} agissent sur l'histoire locale et satisfont
ce contrôle de non-factorisation.

\begin{ablation}[Nécessité des deux feuillets dans le parcours électronique]
\label{abl:hojas-lector-two-way-rev11}
Le lecteur chronologique produit
\[
 (A_D,B_D,C_D)=
 \begin{cases}
 (80,54,6),&\substack{\text{mot chronologique conjoint des deux feuillets ;}\\
 B_D=D_1(\Gamma_e)=54},\\
 (80,84,-8),&\text{projection additive},\\
 (80,60,13),&\text{projection multiplicative}.
 \end{cases}
\]
Le triplet électronique complet exige l'histoire conjointe des deux feuillets
arithmétiques.
\end{ablation}

\subsection{Dénombrements exacts sur les 324 parcours}

Le lecteur chronologique produit quatorze profils. Le tableau suivant en donne le
dénombrement complet ; \(A_D=D_{27}+D_{36}\), \(B_D=D_1\) et
\(C_D=(D_4-D_3)/2\).

\begin{longtable}{@{}rrr r@{}}
\caption{Dénombrement exact des quatorze profils du lecteur chronologique sur les
324 parcours orientés.}
\label{tab:censo-perfiles-kd-rev11}\\
\toprule
\(A_D\)&\(B_D\)&\(C_D\)&multiplicité\\
\midrule
\endfirsthead
\toprule
\(A_D\)&\(B_D\)&\(C_D\)&multiplicité\\
\midrule
\endhead
0&24&0&18\\
40&24&2&36\\
40&54&6&18\\
40&78&27&36\\
40&84&30&36\\
60&78&25&18\\
60&84&28&18\\
80&54&6&18\\
80&54&7&18\\
80&66&2&18\\
80&66&3&36\\
80&66&4&18\\
100&66&0&18\\
100&66&2&18\\
\midrule
\multicolumn{3}{r}{total}&324\\
\bottomrule
\end{longtable}

Le lecteur dodécaphasique produit neuf profils :

\begin{table}[H]
\centering
\caption{Dénombrement exact des neuf profils du lecteur dodécaphasique sur
les 324 parcours orientés.}
\label{tab:censo-perfiles-komega-rev11}
\begin{tabular}{@{}rrr r@{}}
\toprule
\(\Omega\)&\(B_\Omega\)&\(P_6\)&multiplicité\\
\midrule
80&54&6&198\\
56&54&5&68\\
48&54&5&34\\
36&54&4&8\\
20&54&4&4\\
4&54&2&4\\
40&54&4&4\\
24&54&4&2\\
24&54&2&2\\
\midrule
\multicolumn{3}{r}{total}&324\\
\bottomrule
\end{tabular}
\end{table}

Les lecteurs coïncident comme triplets sur (18) parcours, toujours avec la valeur
\((80,54,6)\), et la paire \(\mathbf K\) prend (40) valeurs. Dans la section
canonique \(B1\) à 81 positions, \(K_D\) est constant dans 41 des 56
familles et raffine le quotient en 77 classes ; \(K_\Omega\) est constant dans 50
et produit 62 classes ; la paire en produit 78. Dans le relèvement à 324
orientations, \(K_D\) produit 146 classes et n'est constant dans aucune des
familles ; \(K_\Omega\) en produit 100 et descend dans 17 familles ; la paire en produit
151 et n'est constante dans aucune. Ces dénombrements prouvent que les deux
projections retiennent des mémoires distinctes et fixent le domaine de chaque
affirmation de descente.

\subsection{Traces vectorielles et spectres conjoints des \(13\) multisections}

Dans la section canonique \(B1\), nous notons
\(\ell_j,\nu_j,d_j,u_j\) les multisections de génération \(j\) des
secteurs leptonique chargé, neutrino, quark de type down et quark de type
up, respectivement. Pour chaque lecteur
\(r\in\{D,\Omega\}\) et chaque composante \(a\in\{1,2,3\}\), nous définissons
l'opérateur diagonal
\[
 K_{M,a}^r e_\Gamma=K_{r,a}(\Gamma)e_\Gamma,
 \qquad
 \mathbf K_M^r=(K_{M,1}^r,K_{M,2}^r,K_{M,3}^r).
\]
Les trois opérateurs de \(\mathbf K_M^r\) sont autoadjoints et commutent. Leur
trace normalisée s'entend donc composante par composante,
\[
 \overline{\mathbf K}_M^r
 =\frac{1}{\dim M}
 \bigl(\operatorname{Tr}K_{M,1}^r,
       \operatorname{Tr}K_{M,2}^r,
       \operatorname{Tr}K_{M,3}^r\bigr),
\]
et leur spectre est le spectre conjoint, c'est-à-dire la distribution avec
multiplicité des triplets simultanés. Les treize traces vectorielles sont :

\begin{longtable}{@{}c r c c@{}}
\caption{Traces vectorielles normalisées des lecteurs \(K_D\) et
\(K_\Omega\) dans les treize multisections canoniques.}
\label{tab:trazas-multisecciones-two-way-rev11}\\
\toprule
\(M\)&\(\dim M\)&\(\overline{\mathbf K}_M^D\)&\(\overline{\mathbf K}_M^\Omega\)\\
\midrule
\endfirsthead
\toprule
\(M\)&\(\dim M\)&\(\overline{\mathbf K}_M^D\)&\(\overline{\mathbf K}_M^\Omega\)\\
\midrule
\endhead
\(\ell_0\)&1&\((80,54,6)\)&\((80,54,6)\)\\
\(\ell_2\)&8&\((65,249/4,17/2)\)&\((80,54,6)\)\\
\(\ell_4\)&16&\((225/4,489/8,21/2)\)&\((141/2,54,11/2)\)\\
\(\nu_1\)&2&\((40,51,29/2)\)&\((60,54,5)\)\\
\(\nu_2\)&4&\((45,63,29/2)\)&\((61,54,5)\)\\
\(\nu_3\)&6&\((170/3,55,34/3)\)&\((206/3,54,11/2)\)\\
\(\nu_4\)&8&\((105/2,237/4,95/8)\)&\((143/2,54,45/8)\)\\
\(d_1\)&2&\((40,51,29/2)\)&\((60,54,5)\)\\
\(d_2\)&4&\((45,63,57/4)\)&\((61,54,5)\)\\
\(d_3\)&6&\((170/3,55,23/2)\)&\((218/3,54,17/3)\)\\
\(d_4\)&8&\((105/2,237/4,49/4)\)&\((149/2,54,23/4)\)\\
\(u_1\)&4&\((65,66,9)\)&\((80,54,6)\)\\
\(u_3\)&12&\((205/3,119/2,113/12)\)&\((74,54,23/4)\)\\
\bottomrule
\end{longtable}

La notation \((a,b,c)^{\times n}\) indique la multiplicité \(n\). Le spectre
conjoint complet de \(\mathbf K_D\) est :

\begingroup\footnotesize
\begin{longtable}{@{}c p{.82\textwidth}@{}}
\caption{Spectres conjoints complets du lecteur \(K_D\) dans les treize
multisections canoniques.}
\\
\toprule
\(M\)&\(\operatorname{Spec}_{\rm conj}\mathbf K_M^D\)\\
\midrule
\endhead
\(\ell_0\)&\((80,54,6)^{\times1}\)\\
\(\ell_2\)&\((0,24,0)^{\times1};(40,78,27)^{\times2};(80,54,7)^{\times1};(80,66,2)^{\times1};(80,66,3)^{\times1};(100,66,0)^{\times1};(100,66,2)^{\times1}\)\\
\(\ell_4\)&\((0,24,0)^{\times1}\); \((40,24,2)^{\times2}\); \((40,54,6)^{\times2}\); \((40,84,30)^{\times2}\); \((60,78,25)^{\times3}\); \((80,66,2)^{\times2}\); \((80,66,3)^{\times3}\); \((80,66,4)^{\times1}\)\\
\(\nu_1\)&\((40,24,2)^{\times1};(40,78,27)^{\times1}\)\\
\(\nu_2\)&\((0,24,0)^{\times1};(40,84,30)^{\times1};(60,78,25)^{\times1};(80,66,3)^{\times1}\)\\
\(\nu_3\)&\((40,24,2)^{\times2};(40,78,27)^{\times1};(60,84,28)^{\times1};(80,54,6)^{\times1};(80,66,3)^{\times1}\)\\
\(\nu_4\)&\((0,24,0)^{\times1};(40,24,2)^{\times1};(40,78,27)^{\times2};(40,84,30)^{\times1};(80,54,7)^{\times1};(80,66,2)^{\times1};(100,66,0)^{\times1}\)\\
\(d_1\)&\((40,24,2)^{\times1};(40,78,27)^{\times1}\)\\
\(d_2\)&\((0,24,0)^{\times1};(40,84,30)^{\times1};(60,78,25)^{\times1};(80,66,2)^{\times1}\)\\
\(d_3\)&\((40,24,2)^{\times2};(40,78,27)^{\times1};(60,84,28)^{\times1};(80,54,6)^{\times1};(80,66,4)^{\times1}\)\\
\(d_4\)&\((0,24,0)^{\times1};(40,24,2)^{\times1};(40,78,27)^{\times2};(40,84,30)^{\times1};(80,54,7)^{\times1};(80,66,3)^{\times1};(100,66,2)^{\times1}\)\\
\(u_1\)&\((40,54,6)^{\times1};(60,78,25)^{\times1};(80,66,2)^{\times1};(80,66,3)^{\times1}\)\\
\(u_3\)&\((40,24,2)^{\times2}\); \((40,78,27)^{\times2}\); \((60,84,28)^{\times1}\); \((80,54,6)^{\times3}\); \((80,66,3)^{\times1}\); \((80,66,4)^{\times1}\); \((100,66,0)^{\times1}\); \((100,66,2)^{\times1}\)\\
\bottomrule
\end{longtable}
\label{tab:espectros-kd-multisecciones-rev11}
\endgroup

Le spectre conjoint complet de \(\mathbf K_\Omega\) sur les mêmes fibres est :

\begingroup\footnotesize
\begin{longtable}{@{}c p{.82\textwidth}@{}}
\caption{Spectres conjoints complets du lecteur \(K_\Omega\) dans les treize
multisections canoniques.}
\label{tab:espectros-komega-multisecciones-rev11}\\
\toprule
\(M\)&\(\operatorname{Spec}_{\rm conj}\mathbf K_M^\Omega\)\\
\midrule
\endfirsthead
\toprule
\(M\)&\(\operatorname{Spec}_{\rm conj}\mathbf K_M^\Omega\)\\
\midrule
\endhead
\(\ell_0\)&\((80,54,6)^{\times1}\)\\
\(\ell_2\)&\((80,54,6)^{\times8}\)\\
\(\ell_4\)&\((24,54,2)^{\times1};(56,54,5)^{\times4};(80,54,6)^{\times11}\)\\
\(\nu_1\)&\((40,54,4)^{\times1};(80,54,6)^{\times1}\)\\
\(\nu_2\)&\((4,54,2)^{\times1};(80,54,6)^{\times3}\)\\
\(\nu_3\)&\((36,54,4)^{\times1};(56,54,5)^{\times1};(80,54,6)^{\times4}\)\\
\(\nu_4\)&\((36,54,4)^{\times1};(56,54,5)^{\times1};(80,54,6)^{\times6}\)\\
\(d_1\)&\((40,54,4)^{\times1};(80,54,6)^{\times1}\)\\
\(d_2\)&\((4,54,2)^{\times1};(80,54,6)^{\times3}\)\\
\(d_3\)&\((36,54,4)^{\times1};(80,54,6)^{\times5}\)\\
\(d_4\)&\((36,54,4)^{\times1};(80,54,6)^{\times7}\)\\
\(u_1\)&\((80,54,6)^{\times4}\)\\
\(u_3\)&\((56,54,5)^{\times3};(80,54,6)^{\times9}\)\\
\bottomrule
\end{longtable}
\endgroup

Les traces peuvent coïncider entre multisections de types différents ; le type,
la multiplicité et le spectre complet demeurent des données de l'opérateur.
Le certificat reproduit ces tableaux ligne par ligne.

\HMTClaim{MD-R-MASS-FIBRE-OPERATORS}
\subsection{Opérateurs de fibre et naturalité du changement de base}

Nous distinguons la fibre canonique \(F_M^{B1}\), formée des positions de
la multisection \(M\) avec leur orientation publiée, et la fibre orientée
\(F_M^{\rm or}\), formée de ses quatre relèvements par position. Pour
\(s\in\{B1,{\rm or}\}\), soit
\(\mathcal H_{M,s}=\bigoplus_{\Gamma\in F_M^s}\mathbb C e_\Gamma\). Nous définissons
\begin{equation}
 B_{M,s}^r e_\Gamma=\kappa_r(\Gamma)e_\Gamma,
 \qquad
 R_{\rm act}^{B_{M,s}^r}
 =\exp\bigl[(\log R_{\rm act})B_{M,s}^r\bigr],
 \qquad r\in\{D,\Omega\}.
 \label{eq:operadores-fibra-two-way-rev11}
\end{equation}
Chaque \(B_{M,s}^r\) est autoadjoint et diagonal dans la base des parcours, de sorte que
\(R_{\rm act}^{B_{M,s}^r}\) est positif. Pour une masse de base positive
\(\widehat M_{M,s}^{(0)}\), la forme opératorielle conditionnée à chaque lecteur est
\begin{equation}
 \widehat M_{M,s}^{\beta,r}
 =R_{\rm act}^{B_{M,s}^r/2}\,
  \widehat M_{M,s}^{(0)}\,
  R_{\rm act}^{B_{M,s}^r/2}.
 \label{eq:ley-operatoria-fibra-rev11}
\end{equation}
Cette expression est positive sans imposer la commutativité. Lorsque
\([\widehat M_{M,s}^{(0)},B_{M,s}^r]=0\), elle se réduit à
\(\widehat M_{M,s}^{(0)}R_{\rm act}^{B_{M,s}^r}\).
Un changement de base unitaire \(U\) transforme
\(B_{M,s}^r\mapsto UB_{M,s}^rU^*\) et conserve spectre, trace, déterminant et
polynôme caractéristique. La fibre électronique canonique \(F_e^{B1}\) est de
rang un et se réduit au scalaire du
\cref{thm:reduccion-electron-two-way-rev11}.
Son relèvement orienté est de rang quatre : les parcours \((++)\) et \((--)\)
partagent le profil électronique, tandis que les orientations mixtes conservent
des profils distincts.

Le déterminant normalisé
\[
 \operatorname{ndet}(R_{\rm act}^{B_{M,s}^r})
 =R_{\rm act}^{\operatorname{Tr}(B_{M,s}^r)/\dim F_M^s}
\]
est naturel par permutation de la fibre. L'additivité logarithmique le distingue
comme moyenne géométrique, mais il n'est pas utilisé comme substitut de la réalisation physique.
Le résultat exact du raffinement reste la paire
\((B_{M,s}^D,B_{M,s}^\Omega)\) ; la section scalaire d'une espèce s'obtient
par le morphisme saveur--génération--parcours déjà construit dans
\cref{eq:morfismo-sabor-ruta-cerrado-rev2,eq:realizacion-fisica-sabor-ruta-cerrada-rev2} :
il sélectionne le parcours persistant par des données non métrologiques, puis applique le
caractère de sa signature. Une fonctionnelle scalaire arbitraire
\(S(B_{M,s}^D,B_{M,s}^\Omega)\) ne définit pas cette sélection et ne demeure qu'un
contrôle empêchant la réintroduction de la masse cible.

\subsection{Échelle absolue, correction relative et sensibilités}

La généalogie de la droite d'action s'écrit
\[
 \Sigma_H=\varphi\alpha_{\rm HMT}^{16},
 \qquad \operatorname{ord}_{10}(\Sigma_H)=-34,
\]
\[
 \delta_{2w}=H_5-\eta_{\rm ret}
 =\frac{90}{\pi}\mathscr C_\pi(\alpha_{\rm HMT}),
 \qquad R_{\rm act}=\frac{H_5}{\eta_{\rm ret}},
\]
\[
 \hbar_{\rm pre}=H_5\,10^{-34}\mathcal U_S,
 \qquad
 \hbar_{\rm ret}=\eta_{\rm ret}\,10^{-34}\mathcal U_S.
\]
Le facteur \(10^{-34}\mathcal U_S\) fixe la section absolue et se simplifie dans les
rapports. Pour le lecteur dodécaphasique,
\begin{equation}
 \log\!\left[
 \frac{(m_X^{\beta,\Omega}/m_e^{\beta})}
 {(m_X^{(0)}/m_e^{(0)})}\right]
 =\bigl[(\Omega_X-80)+(P_{6,X}-6)\Delta_4\bigr]
 \log R_{\rm act}.
 \label{eq:razon-masa-omega-rev11}
\end{equation}
Le terme commun \(-54\alpha_{\rm HMT}\) se simplifie également. Pour le lecteur
chronologique,
\begin{equation}
 \log\!\left[
 \frac{(m_X^{\beta,D}/m_e^{\beta})}
 {(m_X^{(0)}/m_e^{(0)})}\right]
 =\bigl[\kappa_D(\Gamma_X)-\kappa_e\bigr]\log R_{\rm act}.
 \label{eq:razon-masa-D-rev11}
\end{equation}
Avec \(\log R_{\rm act}=6.932817628081925\ldots\times10^{-10}\),
\[
 \frac{\partial\log m}{\partial\Omega}=6.9328\times10^{-10},
 \quad
 \frac{\partial\log m}{\partial P_6}
 =\Delta_4\log R_{\rm act}=7.7090\times10^{-14},
\]
\[
 \frac{\partial\log m}{\partial\alpha}
 =-54\log R_{\rm act}=-3.74372\times10^{-8},
 \quad
 \left.\frac{\partial\log m}{\partial\Delta_4}\right|_e
 =6\log R_{\rm act}=4.15969\times10^{-9}.
\]
Ces dérivées sont des sensibilités du modèle et restent distinctes de
l'incertitude expérimentale d'une masse de comparaison.

\subsection{Appariement physique et secteur nul}

L'appariement typé suit la chaîne
\[
 \begin{aligned}
 &\text{registre physique typé}
 \longrightarrow\text{multisection ou représentation}
 \longrightarrow\text{fibre de routes}\longrightarrow\text{registre}\\
 &\longrightarrow\text{signature et caractère}
 \longrightarrow(B_{M,s}^D,B_{M,s}^\Omega)
 \longrightarrow\text{carte dimensionnelle}
 \longrightarrow\text{confrontation externe}.
 \end{aligned}
\]
L'inventaire physique typé distingue les enregistrements par objet, représentation et codomaine ; les 324 parcours constituent
le moteur interne ; les 471 lignes du catalogue PDG sont des propriétés externes.
Ces cardinaux appartiennent à des domaines différents. La flèche est construite
à partir des descripteurs physiques non métrologiques, de la multisection et du parcours ; une
recherche par proximité avec une masse relève du régime d'étalonnage.

Le secteur nul bifurque avant le caractère positif. Pour le photon et les gluons,
\[
 \kappa_{\rm null}=\mathrm{N/A},
 \qquad m_{\rm null}=0.
\]
L'expression \(R_{\rm act}^{0}=1\) est l'unité du caractère positif. Les
composés concatènent les
enregistrements avant d'évaluer la signature ; les secteurs topologiques conservent fusion et
tresse ; une section virtuelle est classée selon son horizon fini de
prolongement.

\subsection{Indépendance du constructeur de parcours vis-à-vis des masses et clôture de l'opérateur}

La vérification structurelle utilise exclusivement le dénombrement des parcours,
l'enregistrement \(81\times108\) et le transducteur des 81 positions. Elle vérifie leur
couplage ligne par ligne, la couverture
\(81/56/13/324\), les naturalités, le triplet électronique \((80,54,6)\), les
dénombrements (14) et (9), la coïncidence (18/324), les
(40) paires, les descentes différenciées dans \(B1\) et dans les 324 orientations,
la périodicité 54, l'inversion par décalage de 27 et les ablations de feuillet. Le
secteur nul appartient au domaine extérieur au caractère positif. Le calcul est
effectué avant l'introduction de masses, de CODATA, de PDG, de résidus métrologiques ou
d'étiquettes nominales.

Ce certificat n'évalue ni \(\alpha_{\rm HMT}\), ni \(\Delta_4\),
ni \(R_{\rm act}\), ni \(\kappa_e\), ni la carte MeV. Ces évaluations relèvent de
la chaîne arithmétique de la droite d'action et de la réalisation énergétique du
\cref{thm:reduccion-electron-two-way-rev11} ; cette chaîne arithmétique demeure
distincte de la vérification combinatoire des lecteurs.

La comparaison métrologique applique ensuite les opérateurs à l'inventaire typé,
déclare unité, schéma, échelle, date et incertitude, puis effectue des permutations
et une validation par mise en réserve de multisections complètes. Le résultat physique est
régi par le domaine \(81/56/13/324\), par les deux lecteurs de parcours et par
les opérateurs de fibre.

\begin{statusbox}[title={Domaine exact et réalisation physique achevée}]
L'atlas, la chaîne \(C_{81}\to\mathcal L_{108}\) de l'enregistrement chronologique
enrichi de \(108\) itérations,
la droite d'action,
les lecteurs \(K_D,K_\Omega\), leur réduction électronique et les opérateurs de
fibre sont des résultats internes exacts. La paire
\(\mathbf K=(K_D,K_\Omega)\) est le résultat complet sur chaque parcours et
détermine les deux opérateurs diagonaux de toute fibre. Dans la fibre
électronique, les deux lecteurs coïncident. Dans les fibres non centrales, la paire et son
spectre préservent le raffinement complet, tandis que
\(\mathcal S_{\rm phys}\) relie branche, génération, signature fine et, le cas
échéant, orbite de couleur au parcours persistant. Le caractère de ce parcours
donne la coordonnée scalaire avant la confrontation métrologique. Il n'existe donc pas
de disjonction entre opérateur achevé et masse en attente : ce sont deux lectures
typées de la même généalogie. Les deux lecteurs sont codirecteurs, et aucune
proximité avec une masse externe n'intervient dans la sélection.
\end{statusbox}
\HMTClaim{MD-R-MASS-PHYSICAL-SELECTION-CLOSED}

\endgroup
\begingroup
\let\chapter\section
\let\section\subsection
\let\subsection\subsubsection
\let\subsubsection\paragraph
\let\clearpage\relax
% Vista editorial exacta de corpus/source/masas/04_persistencia.tex, líneas 365--590. Las líneas 1--364 ya residen exactamente en 73_rev2_electron_lectores_torres.tex.
\subsection{Persistance de l'état enrichi sous l'extension survivante et conservation de la mémoire de parcours}
Soit \(x\) un état APP, \(\tau\) son mot tritique et \(B_K\) le bloc de mémoire au niveau \(K\). L'état pertinent pour la masse prend la forme

\[
X_K=(x_K,\tau_K,\varepsilon_K,h_K,v_K,c_K,B_K,g_K),
\]
où \(\varepsilon_K\) est l'orientation du feuillet, \((h_K,v_K)\) l'orientation cellulaire, \(c_K\) la retenue et \(g_K\) la phase nonadique. Une extension est survivante lorsqu'elle forme une section compatible des cylindres successifs. Si \(C_{K,j}\), \(0\le j<729\), sont les descendants du cylindre au niveau suivant, la règle d'élagage satisfait

\[
\sum_{j=0}^{728}G_9(C_{K,j})=1.
\]
L'indice unique \(j_K\) produit la transition

\[
P_{K+1}=729P_K+j_K.
\]
Après neuf niveaux, la phase se conserve,

\[
g(K+9)=g(K),
\]
mais le bloc, le préfixe et la mémoire changent. Le retour de phase n'est donc pas un retour d'état. Une particule stable se définit par la compatibilité de la section à travers ces retours, et non par la permanence immobile d'une coordonnée visible.

Le parcours observable \(\gamma_{\rm obs}\) est l'ombre chronologique de l'extension. Le point qui parcourt un diagramme est un curseur de lecture ; la particule est la classe complète de persistance. Cette distinction permet d'interpréter correctement le mouvement quantique : il n'est pas nécessaire de supposer une petite bille suivant une trajectoire classique, mais une histoire cohérente admettant plusieurs projections compatibles.

\Needspace{7\baselineskip}
\subsection{Enregistrement dodécaphasique et signature}
La roue temporelle contient douze événements orientés. Sa décomposition causale est

\[
12=1+5+6:
\]
un événement fixe l'ancrage, cinq déterminent l'incidence visible et six forment les paires opposées qui survivent au quotient orienté. Si \(s_C^{\rm tot}\) est la somme contre-angulaire entière sur les douze événements, l'enregistrement conserve

\[
s_C:=s_C^{\rm tot}\in\mathbb Z,
\qquad
s_C^{(12)}:=\frac{s_C^{\rm tot}}6.
\]
Le premier objet est la coordonnée entière de la signature ; le second, sa lecture angulaire dodécaphasique. Le diviseur n'est pas une normalisation introduite à partir d'une unité externe. C'est l'indice qui apparaît lorsque l'on reconstruit les six paires opposées. La matrice d'incidence dodécaphasique possède une forme normale de Smith de facteurs \((1,12)\) ; le caractère est bien défini sur le quotient parce que le résidu d'orientation se mesure dans cette structure.

Telle est la convention directrice de tout le traité. Les expressions historiques conservées qui désignaient par \(s_C\) la coordonnée déjà divisée se lisent comme \(s_C^{(12)}\) ; leurs valeurs numériques et leurs témoignages demeurent intacts, mais aucune formule en vigueur ne divise deux fois la même coordonnée.

La paire d'enroulements directs est

\[
W_+=6n_A+s_C^{\rm tot},
\qquad
W_-=6n_A-s_C^{\rm tot}.
\]
Puisque \(s_C=s_C^{\rm tot}\) dans la signature publiée, cette formule coïncide avec \(W_\pm=6n_A\pm s_C\). L'exposant contre-angulaire utilise une seule fois \(s_C^{(12)}=s_C/6\) ; il ne divise pas de nouveau une coordonnée déjà normalisée.

L'enregistrement complet donne la signature

\[
\sigma(\Gamma)=
(n_A,s_C,k_{\Delta_4},b_{120},b_\zeta)\in\mathbb Z^5,
\qquad
\zeta_{270}=135\Delta_4^2.
\]
Les cinq composantes sont de types distincts. \(n_A\) compte le développement angulaire direct ; \(s_C\) enregistre la somme contre-angulaire entière ; \(s_C^{(12)}=s_C/6\) est sa lecture angulaire ; \(k_{\Delta_4}\) lit le raffinement global ; \(b_{120}\) marque le canal harmonique de hauteur 120 ; \(b_\zeta\) appartient au correcteur quadratique de hauteur 270. En particulier, \(s_{270}\) et \(\zeta_{270}\) ne sont pas le même objet : le premier est une trace impaire de l'opérateur bicouche ; le second est une coordonnée quadratique de l'enregistrement.

Le mot électronique

\[
\tau_{12}=+++---+++---
\]
a une somme linéaire nulle, mais sa signature paire n'est pas nulle :

\[
(Q_3,Q_4,Q_6,\rho_0,w)=(-12,-4,12,-6,12).
\]
Il ne faut pas non plus l'identifier à la signature brute du cycle électronique \((24,0,12,0,-12)\), ni à l'origine affine \(v_e=0\). Il s'agit de quatre niveaux de lecture distincts : mot, signature paire, signature brute et origine du caractère relatif.

\Needspace{7\baselineskip}
\subsection{Angle, contre-angle et caractère}
Pour rendre la construction autonome, nous fixons ici les trois coordonnées internes qui interviennent dans toute la loi :

\[
\alpha_{\rm HMT}
=\frac{7297352569283800997285105472380663}{10^{36}},
\]
\[
\Delta_4
=\frac{\pi_{\rm HMT}-e_{\rm HMT}\log\pi_{\rm HMT}}{270}
\left(1+\frac{\pi_{\rm HMT}}{729}\right),
\]
et, pour \(a>0\) et \(\phi>0\),

\[
H_5(a,\phi)
=\frac12\sqrt{\frac{1000a}{\phi}}-a
+\frac{9a^2}{16}-\frac{5a^3}{9}
+\frac{7a^4}{48}-\frac{a^5}{54}.
\]
Dans la suite,

\[
H_5:=H_5(\alpha_{\rm HMT},\varphi_{\rm HMT}).
\]
Ces définitions n'importent aucune constante métrologique externe : \(\pi_{\rm HMT},e_{\rm HMT},\varphi_{\rm HMT}\) sont les évaluations coinductives du lecteur HMT et \(\alpha_{\rm HMT}\) est le résultat rationnel de la clôture dodécaphasique. La constante interne \(\alpha_{\rm HMT}\) précède le contre-angle. La branche régionale détermine d'abord

\[
A=1000\alpha_{\rm HMT},
\qquad
D_A=\exp\!\left(-\frac{100\pi\alpha_{\rm HMT}}9\right),
\]
\[
\mathcal C_\pi(\alpha_{\rm HMT})
=169\alpha_{\rm HMT}^{6}
+\frac{D_A\alpha_{\rm HMT}^{7}}
       {1-D_A\alpha_{\rm HMT}},
\qquad
y_*=\frac{\pi}{90}(H_5+\alpha_{\rm HMT})
-\mathcal C_\pi(\alpha_{\rm HMT}),
\]
\[
C^*=\frac{180}{\pi}y_*.
\]
Ici, \(A\) et \(C^*\) sont des coordonnées angulaires déjà construites, et non des grandeurs d'action. Si

\[
\alpha_\angle:=\alpha_{\rm HMT}\,{}^\circ,
\qquad
\eta_{\rm ret}^{(\deg)}:=\frac{C^*}{2}-\alpha_\angle,
\]
alors

\[
C^*=2\bigl(\eta_{\rm ret}^{(\deg)}+\alpha_\angle\bigr)
\]
est l'identité inverse de retour, et non une seconde règle de sélection du contre-angle. Cette écriture évite de confondre la mantisse angulaire \(\eta_{\rm ret}^{(\deg)}\) avec la grandeur dimensionnelle \(\hbar_{\rm ret}\). Les cartes en radians s'obtiennent par multiplication par \(\pi/180\). Les canaux conjugués sont

\[
q_\pm=\exp\left[-\frac\pi{180}(A\pm C^*)\right].
\]
Le caractère central d'une signature est

\[
\chi_0(\sigma)=
\exp\left(
n_AA_{\rm rad}
+\frac{s_C}{6}C^*_{\rm rad}
+k_{\Delta_4}\Delta_4
+b_\zeta\zeta_{270}
\right).
\]
La coordonnée \(b_{120}\) demeure dans la signature, mais sa contribution est incorporée dans la famille de tours. Cette convention empêche de compter deux fois la même hauteur.

Le caractère est adimensionnel. Sa fonction n'est pas de devenir une constante d'action, mais d'agir par homothéties positives sur une droite de masse déjà typée. Pour une paire d'états \(X,Y\),

\[
\frac{m_Y^{(0)}}{m_X^{(0)}}
=\chi_{\rm full}(\sigma_Y-\sigma_X,\eta_Y-\eta_X).
\]
Cette formule sépare les rapports de base des masses du choix commun de carte absolue. Le quotient bêta complet incorpore ensuite \(R_{\rm act}^{\kappa_Y-\kappa_X}\).

\Needspace{7\baselineskip}
\subsection{Semi-groupe global des tours de masse et développement harmonique convergent}
Soit \(R^2=I\), avec les projecteurs

\[
P_\pm=\frac{I\pm R}{2},
\qquad
x=\frac{\pi A}{180},
\qquad
y=\frac{\pi C^*}{180},
\]
et

\[
T=e^{-xI-yR}=q_+P_++q_-P_-.
\]
Les traces paire et impaire sont

\[
c_n=\operatorname{Tr}(T^n)=q_-^n+q_+^n
=2e^{-nx}\cosh(ny),
\]
\[
s_n=-\operatorname{Tr}(RT^n)=q_-^n-q_+^n
=2e^{-nx}\sinh(ny).
\]
Elles satisfont

\[
c_n^2-s_n^2=4e^{-2nx},
\]
\[
c_{m+n}=\frac{c_mc_n+s_ms_n}{2},
\qquad
s_{m+n}=\frac{s_mc_n+c_ms_n}{2},
\]
et la récurrence linéaire d'ordre deux

\[
X_{n+2}=c_1X_{n+1}-e^{-2x}X_n,
\qquad X\in\{c,s\},
\qquad c_0=2,\quad s_0=0.
\]
L'ensemble des hauteurs ne se réduit pas à huit corrections visibles. C'est le semi-groupe

\[
\langle90,120\rangle
=\{90,120\}\cup\{30r:r\ge6\}.
\]
Les hauteurs \(90,120,180,210,240,270,360,420\) sont les premiers témoins généalogiques utiles, et non le domaine entier. Le caractère raffiné est

\[
\log\chi_{\rm full}(\sigma,\eta)
=\log\chi_0(\sigma)
+\sum_{h\in\langle90,120\rangle}N_h(\sigma,\eta)s_h,
\]
avec

\[
N_h(\sigma,\eta)=\eta_h+b_{120}\delta_{h,120}.
\]
Ainsi, \(s_{120}\) apparaît une seule fois, et à la hauteur 270, \(N_{270}s_{270}\) et \(b_\zeta\zeta_{270}\) restent distincts. La convergence absolue exige

\[
\sum_h |N_hs_h|<\infty.
\]
La complétude du semi-groupe garantit la capacité de représentation des rapports positifs ; elle ne choisit pas, à elle seule, les coefficients physiques d'une espèce. La règle prospective de sélection, lorsqu'elle intervient, doit procéder du parcours, de l'enregistrement généalogique, de la retenue, de la mémoire et de l'enroulement avant l'introduction de la masse externe.

\Needspace{7\baselineskip}

\endgroup
\begingroup
\let\chapter\section
\let\section\subsection
\let\subsection\subsubsection
\let\subsubsection\paragraph
\let\clearpage\relax
% Vista editorial exacta de corpus/source/masas/06_operadores.tex, líneas 50--113. Las líneas 1--29 y 30--49 ya residen exactamente en 73_rev2_electron_lectores_torres.tex.
\subsection{Lecteurs bidirectionnels \(K_D\) et \(K_\Omega\) : profils, coïncidences et opérateurs par multisection}
Le parcours orienté produit deux familles d'invariants. Pour les défauts temporels \(D_j\) et la mémoire de retour \(\Omega\), nous définissons

\[
K_D=\left(D_{27}+D_{36},D_1,\frac{D_4-D_3}{2}\right),
\qquad
K_\Omega=(\Omega,54,P_6).
\]
Sur les 324 parcours apparaissent 14 profils distincts de \(K_D\), 9 profils de \(K_\Omega\), 40 paires conjointes et 18 coïncidences. L'unique profil commun du parcours électronique est

\[
K_D(\Gamma_e)=K_\Omega(\Gamma_e)=(80,54,6).
\]
Ces lecteurs ne sont pas des étiquettes décoratives. \(K_D\) mesure les écarts du mot ; \(K_\Omega\) mesure le retour enrichi. Leur égalité au centre certifie la clôture bidirectionnelle de rang un. Hors du centre, ils produisent des opérateurs de fibre dont les spectres complets doivent être présentés, même lorsqu'une réalisation physique ultérieure sélectionne une section scalaire.

\Needspace{7\baselineskip}
\subsection{Raffinement harmonique propre à chaque état}
Le raffinement complet se situe dans l'espace de coefficients

\[
\mathcal E_{90,120}
=\left\{\eta=(\eta_h)_{h\in\langle90,120\rangle}:
\sum_h|\eta_hs_h|<\infty\right\}.
\]
L'application recherchée ne reçoit pas une masse. Elle reçoit le parcours enrichi :

\[
\mathfrak T:\Gamma_{\rm enr}\longrightarrow\mathcal E_{90,120},
\qquad
\Gamma\longmapsto\eta(\Gamma).
\]
Ses coordonnées doivent être des fonctions explicites des retours, de la retenue, de l'enroulement, de la mémoire et de la frontière. La condition de naturalité est

\[
\eta(r_{n,m}\Gamma)=r^{\mathcal E}_{n,m}\eta(\Gamma),
\]
et la convergence doit s'accompagner d'une borne de reste calculable. Le caractère complet agit alors par

\[
\widehat M_F
=\sum_{\gamma\in F}
m_{\rm clk}\,
\chi_{\rm full}(\sigma(\gamma),\eta(\gamma))
|\gamma\rangle\langle\gamma|.
\]
La capacité de reconstruire une différence numérique donnée démontre la complétude représentative de la hiérarchie ; elle ne démontre pas, à elle seule, que \(\mathfrak T\) a choisi ses coefficients. La différence entre les deux affirmations est préservée dans chaque ligne quantitative.

\Needspace{7\baselineskip}
\subsection{Scalarisation par couplage}
Soit \(P_{X,a}\) le projecteur déterminé par l'état \(X\) et l'appareil ou le canal physique \(a\). La distribution de masse est

\[
\mu_{X,a}(B)
=\langle\Psi_X,
P_{X,a}E_{\widehat M_F}(B)P_{X,a}\Psi_X\rangle.
\]
Il existe une droite scalaire lorsque la compression satisfait

\[
P_{X,a}\widehat M_FP_{X,a}=m_XP_{X,a}.
\]
Tel est le sens précis de « mesurer, c'est se coupler ». La valeur numérique n'était pas cachée dans une cellule nominale ; elle apparaît comme valeur propre d'une compression physique d'un opérateur déjà construit. Si la compression conserve plusieurs valeurs propres, la prédiction correcte est un multiplet ou une distribution, et non une moyenne choisie pour coïncider avec une référence.

\Needspace{7\baselineskip}


\endgroup
\begingroup
\let\chapter\section
\let\section\subsection
\let\subsection\subsubsection
\let\subsubsection\paragraph
\let\clearpage\relax
\subsection{Sélection physique comme relation naturelle}
En dehors du centre électronique, une espèce ne doit pas choisir par décret une cellule unique. Les douze multisections non centrales sont des fibres stables sous l’inversion cellulaire et n’admettent généralement pas de section ponctuelle équivariante. La réalisation correcte commence par une relation

\[
\mathscr S\subset \mathcal P\times\mathcal R_{324},
\qquad
\mathscr S(X)=\{\gamma:(X,\gamma)\in\mathscr S\},
\]
dont la valeur peut être un point, une orbite ou une fibre. Pour l’électron, \(\mathscr S(e)=\{\gamma_e\}\) est de rang un parce que \((5,5)\) est l’unique point fixe. Pour la couleur, la saveur et la génération, \(\mathscr S(X)\) conserve l’orbite correspondante jusqu’à ce qu’un couplage physique la comprime.

La relation est admissible si elle satisfait simultanément les conditions suivantes :

\begin{enumerate}
\item elle dépend uniquement de la représentation, de l’orientation, de la mémoire, de la frontière, de la charge, de la saveur, de la couleur et de la génération construites avant la comparaison ;

\item elle commute avec la conjugaison de particule et d’antiparticule ;
\item elle est équivariante relativement à l’action chromatique et à l’inversion cellulaire ;
\item elle est compatible avec le raffinement, le retour nonadique et la lecture dodécaphasique ;
\item elle conserve toute multiplicité qui n’a pas été éliminée par un couplage ;
\item elle reste invariante lorsqu’on permute ou remplace les valeurs externes de masse.
\end{enumerate}
Le sixième point est le contrôle décisif. Si la modification de la colonne expérimentale change \(\mathscr S(X)\), la construction est rétrospective. Si la relation demeure et que seul le résidu de comparaison change, la direction causale est préservée.

\Needspace{7\baselineskip}

\endgroup

\section{Morphisme de scalarisation et obstructions de la réalisation dimensionnelle}
Hors de la fibre électronique de rang un, le spectre ne sélectionne pas à lui seul une valeur propre physique. Un scalaire exige un projecteur, un état ou une fonctionnelle de couplage déclarés avant la confrontation externe.
